\documentclass[11pt,a4paper]{article}
\usepackage[margin=1.8cm]{geometry}
\usepackage{amsmath,amssymb,array,tabularx,fontspec,cancel,tikz,graphicx}
\setmainfont{FreeSerif}
\newfontfamily\urdufont[Script=Arabic]{FreeSerif}
\TeXXeTstate=1
\newcommand{\urdu}[1]{{\urdufont\beginR #1\endR}}
\newcommand{\lr}[1]{\beginL #1\endL}
\setlength{\parindent}{0pt}
\setlength{\parskip}{0pt}
\renewcommand{\arraystretch}{1.15}
\newcommand{\ut}[1]{\text{{\urdufont\beginR #1\endR}}}
% \uq{N}{Urdu text}: question with its printed number on the right
\newcommand{\uq}[2]{\par\medskip{\urdufont\raggedleft\beginR\textbf{\lr{#1.}}\ #2\endR\par}}
% continuation paragraph (no number)
\newcommand{\usc}[1]{\par\smallskip{\urdufont\raggedleft\beginR #1\endR\par}}
% centred math line
\newcommand{\mline}[1]{\par\smallskip\centerline{\lr{$\displaystyle #1$}}}
\newcommand{\tfig}[1]{\par\smallskip\centerline{{\TeXXeTstate=0 \begin{tikzpicture}[scale=0.8,font=\small]#1\end{tikzpicture}}}\par}
\newcommand{\pfig}[2]{\par\smallskip\centerline{\includegraphics[width=#2]{#1}}\par}
% data table (LTR table inside a centred line)
\newcommand{\dtab}[2]{\par\smallskip\centerline{\lr{\begin{tabular}{#1}\hline #2\\ \hline\end{tabular}}}\par\smallskip}
\newcommand{\nums}[1]{\par\smallskip\centerline{\lr{\parbox{0.9\linewidth}{#1}}}\par\smallskip}
\newcommand{\foodsurvey}{ایک سروے کیا گیا جس میں \lr{130} لوگوں سے ان کے پسندیدہ کھانوں کے بارے میں پوچھا گیا۔ سروے کے نتائج نے درج ذیل اعداد و شمار ظاہر کیے۔ \lr{40} لوگوں نے کہا کہ وہ نہاری کو پسند کرتے ہیں۔ \lr{65} لوگوں نے کہا کہ انہیں بریانی پسند ہے۔ \lr{50} لوگوں نے کہا کہ انہیں قورمہ پسند ہے۔ \lr{20} لوگوں نے کہا کہ انہیں نہاری اور بریانی پسند ہے۔ \lr{35} لوگوں نے کہا کہ انہیں بریانی اور قورمہ پسند ہے۔ \lr{27} لوگوں نے کہا کہ انہیں نہاری اور قورمہ پسند ہے۔ \lr{12} لوگوں نے کہا کہ انہیں تینوں کھانے نہاری، بریانی اور قورمہ پسند ہیں۔}
\newcommand{\mallemp}{ایک شاپنگ مال میں \lr{1} سے \lr{150} لیبل والے \lr{150} ملازمین ہیں جو یونیورسل سیٹ \lr{U} کو ظاہر کرتے ہیں۔ ملازمین درج ذیل زمروں میں آتے ہیں: سیٹ \lr{A: 40} ملازمین جن کی تنخواہ کی حد \lr{30} ہزار تا \lr{45} ہزار ہے۔ جس پر \lr{50} سے \lr{89} تک کا لیبل لگایا گیا ہے۔ سیٹ \lr{B: 50} ملازمین جن کی تنخواہ کی حد \lr{50} ہزار تا \lr{80} ہزار ہے۔ جس کا لیبل \lr{101} سے \lr{150} تک ہے۔ سیٹ \lr{C: 60} ملازمین جن کی تنخواہ کی حد \lr{100} ہزار تا \lr{150} ہزار ہے۔ \lr{1} سے \lr{49} اور \lr{90} سے \lr{100} تک کا لیبل لگا ہوا ہے۔}
\newcommand{\sportsch}{ایک سیکنڈری سکول میں \lr{125} طلبہ میں سے ہر ایک درج ذیل میں سے کم از کم ایک کھیل میں حصہ لیتا ہے۔ کرکٹ، فٹ بال یا ہاکی۔ \lr{60} طلبہ کرکٹ کھیلتے ہیں۔ \lr{70} طلبہ فٹ بال کھیلتے ہیں۔ \lr{40} طلبہ ہاکی کھیلتے ہیں۔ \lr{25} طلبہ کرکٹ اور فٹ بال دونوں کھیلتے ہیں۔ \lr{15} طلبہ فٹ بال اور ہاکی دونوں کھیلتے ہیں۔ \lr{10} طلبہ کرکٹ اور ہاکی دونوں کھیلتے ہیں۔}
\newcommand{\itvenn}{ملحقہ وین شکل پر غور کریں جس میں ایک آئی ٹی ادارے میں مختلف کورسز میں داخلہ لینے والے طلبہ کو ظاہر کیا گیا ہے۔ \lr{U=} آئی ٹی ادارے میں داخلہ لینے والے طلبہ، \lr{A=} اپلائیڈ روبوٹکس میں داخلہ لینے والے طلبہ، \lr{G=} گیم ڈویلپمنٹ میں داخلہ لینے والے طلبہ \lr{W=} ویب ڈیزائننگ میں داخلہ لینے والے طلبہ۔}
\newcommand{\shopsurvey}{ایک شاپنگ مال میں \lr{130} صارفین کا سروے کیا گیا۔ جس میں ان سے خریداری کی ترجیحات کے بارے میں پوچھا گیا۔ سروے کے نتائج نے درج ذیل اعداد و شمار ظاہر کیے۔ \lr{57} صارفین نے گارمنٹس کی خریداری کی۔ \lr{50} صارفین نے سنگھار کا سامان خریدا۔ \lr{46} صارفین نے الیکٹرونکس کی اشیا خریدی۔ \lr{31} صارفین نے گارمنٹس اور سنگھار دونوں کا سامان خریدا۔ \lr{25} صارفین نے گارمنٹس اور الیکٹرانکس دونوں کی خریداری کی۔ \lr{21} صارفین نے سنگھار اور الیکٹرونکس دونوں کی خریداری کی۔ \lr{12} صارفین نے تینوں مصنوعات یعنی گارمنٹس، سنگھار اور الیکٹرونکس کی خریداری کی۔}
\newcommand{\fzA}[2]{\uq{#1}{تجزی کریں: \lr{$#2$}}}
\newcommand{\fzB}[2]{\uq{#1}{درج ذیل جملے کی تجزی کریں۔ \lr{$#2$}}}
\newcommand{\abprop}{ذیل میں دیے گئے سیٹ \lr{A,B} اور \lr{C} کے لیے خصوصیات کو ثابت کریں: (الف) یونین کی خاصیت تلازم (ب) تقاطع کی خاصیت تلازم (ج) یونین کی تقاطع پر خاصیت تقسیمی (د) تقاطع کی یونین پر خاصیت تقسیمی}
\newcommand{\examtxt}{سالانہ امتحان میں شامل ہونے والے \lr{180} طلبہ میں سے \lr{120} نے ریاضی کا امتحان پاس کیا، \lr{90} نے سائنس کا امتحان پاس کیا۔ اور \lr{60} نے ریاضی اور سائنس دونوں ٹیسٹ پاس کیے۔}
\newcommand{\hcfdiv}[2]{\uq{#1}{درج ذیل جملے کا عاد اعظم بذریعہ تقسیم معلوم کریں:}\mline{#2}}
\newcommand{\hcffac}[2]{\uq{#1}{عاد اعظم بذریعہ تجزی معلوم کریں: \lr{$#2$}}}
\newcommand{\sqrdiv}[2]{\uq{#1}{درج ذیل جملے کا جذر المربع بذریعہ تقسیم معلوم کریں:}\mline{#2}}
\newcommand{\shade}[3]{\uq{#1}{درج ذیل یک درجی غیر مساواتوں کے حل کا خطہ سایہ دار کر کے ظاہر کریں: \lr{$#2,\ #3$}}}
\newcommand{\trig}[2]{\uq{#1}{درج ذیل تکونیاتی اکائی کو ثابت کریں:}\mline{#2}}
\newcommand{\solvetri}[2]{\uq{#1}{درج ذیل مثلث کو حل کریں جب کہ: \lr{$#2$}}}
\newcommand{\quarter}[2]{\uq{#1}{اگر \lr{$\sin\frac{\pi}{4}$} اور \lr{$\cos\frac{\pi}{4}$} کی قیمت \lr{$\frac{1}{\sqrt{2}}$} ہو تو درج ذیل مقدار معلوم کریں: \lr{$#2$}}}
\newcommand{\trigb}[2]{\uq{#1}{درج ذیل تکونی اکائی کو ثابت کریں:}\mline{#2}}
\newcommand{\ptsq}[2]{\uq{#1}{نقاط کے دو جوڑے دیے گئے ہیں۔ معلوم کریں کہ آیان نقاط سے بنے والے دو خطوط (الف) متوازی ہیں (ب) عمودی ہیں (ج) کچھ بھی نہیں ہیں \lr{$#2$}}}
\newcommand{\lineform}[2]{\uq{#1}{درج ذیل مساوات کو ڈھلوان قاطع شکل، دو قاطع شکل اور عام شکل میں تبدیل کریں۔}\mline{#2}}
\newcommand{\parperp}[2]{\uq{#1}{درج ذیل میں چیک کریں کہ آیا دو خطوط: (الف) متوازی ہیں (ب) عمودی ہیں (ج) نامتوازی ہیں اور نہ عمودی ہیں۔}\mline{#2}}
\newcommand{\dproof}[2]{\uq{#1}{درج ذیل الجبری جملے کے لیے سادہ استخراجی ثبوتوں کو وضع کریں اور ثابت کریں کہ بائیں طرف برابر ہے دائیں طرف کے: \lr{$#2$}}}
\newcommand{\jproof}[2]{\uq{#1}{ہر مرحلے کی وضاحت کرتے ہوئے درج ذیل کو ثابت کریں: \lr{$#2$}}}
\newcommand{\ptcoord}[2]{\uq{#1}{#2}}
\newcommand{\plotq}[2]{\uq{#1}{درج ذیل دو درجی اور سہ درجی تفاعل کا گراف بنائیں:}\mline{#2}}
\newcommand{\bisect}[2]{\uq{#1}{تصدیق کریں کہ درج ذیل پیمائشوں سے \lr{$\triangle ABC$} کے زاویوں کے ناصف ہم نقطہ ہوتے ہیں:}\mline{#2}}
\newcommand{\perpbis}[2]{\uq{#1}{دی گئی پیمائشوں سے \lr{$\triangle ABC$} بنائیں اور تصدیق کریں کہ مثلث کے عمودی ناصف ہم نقطہ ہوتے ہیں:}\mline{#2}}
\newcommand{\medq}[2]{\uq{#1}{درج ذیل پیمائشوں سے \lr{$\triangle LMN$} بنائیں اور تصدیق کریں کہ مثلث کے وسطانیے ہم نقطہ ہوتے ہیں:}\mline{#2}}

\begin{document}
\begin{center}
  {\Huge\bfseries The Hope Science Academy}\\[3pt]
  {\small Excellence is our standard not goal}\\[2pt]
  {\small 247-E-1, (Opposite LDA School) Johar Town, Lahore. PH: 0300-8194789, 0322-4157001}
\end{center}
\medskip\noindent
\begin{tabularx}{\linewidth}{|X|X|X|X|}
\hline
\textbf{Student Name:} & \textbf{Roll Num:} & \textbf{Class:} 10th & \textbf{Paper Code:} \\ \hline
\textbf{Subject Name:} Mathematics. & \textbf{Time Allowed:} & \textbf{Total Marks:} 2318 & \textbf{Exam Date:} 09-Oct-2026 \\ \hline
\multicolumn{4}{|l|}{\textbf{Exam Syllabus:} FULL BOOK} \\ \hline
\end{tabularx}
\bigskip
\medskip\noindent\textbf{1$\times$445=445}\hfill{\urdufont\beginR\textbf{سوال نمبر \lr{6}۔ مندرجہ ذیل سوالات کے تفصیلی جوابات تحریر کریں۔}\endR}\par
\uq{1}{\lr{x} اور \lr{y} کی قیمت معلوم کیجیے: \lr{$(3x+2)-(4-y)i=5+3i$}}
\uq{2}{\lr{x} اور \lr{y} کی قیمت معلوم کیجیے: \lr{$(2+i)x+(1-2i)y=3+4i$}}
\uq{3}{\lr{x} اور \lr{y} کی قیمت معلوم کیجیے: \lr{$(1-i)x+(2+i)y=4-i$}}
\uq{4}{\lr{x} اور \lr{y} کی قیمت معلوم کیجیے: \lr{$(3x-1)+(2y-3)i=8+7i$}}
\uq{5}{اگر \lr{$z_{1}=2+5i,\ z_{2}=1-3i$} اور \lr{$z_{3}=2+i$} ہو تو تصدیق کیجیے:}
\mline{(z_{1}+z_{2})+z_{3}=z_{1}+(z_{2}+z_{3})}
\uq{6}{اگر \lr{$z_{1}=2+5i,\ z_{2}=1-3i$} اور \lr{$z_{3}=2+i$} ہو تو تصدیق کیجیے:}
\mline{(z_{1}z_{2})z_{3}=z_{1}(z_{2}z_{3})}
\uq{7}{اگر \lr{$\dfrac{(1+i)^{2}}{2-i}=x+iy$} ہو تو \lr{x} اور \lr{y} کی قیمت معلوم کیجیے۔}
\uq{8}{اگر \lr{$(2x+iy)(1-i)=4+2i$} ہو تو \lr{x} اور \lr{y} کی قیمتیں معلوم کیجیے۔}
\uq{9}{\lr{$\dfrac{3+4i}{5-7i}$} کو \lr{$x+iy$} یا \lr{$x+yi$} کی شکل میں ظاہر کیجیے۔}
\uq{10}{اگر \lr{$z_{1}=2+7i$} اور \lr{$z_{2}=4-3i$} ہو تو تصدیق کیجیے کہ:}
\mline{\overline{z_{1}+z_{2}}=\overline{z_{1}}+\overline{z_{2}}}
\uq{11}{اگر \lr{$z_{1}=2+7i$} اور \lr{$z_{2}=4-3i$} ہو تو تصدیق کیجیے کہ:}
\mline{\overline{z_{1}z_{2}}=\overline{z_{1}}\ \overline{z_{2}}}
\uq{12}{اگر \lr{$z_{1}=2+7i$} اور \lr{$z_{2}=4-3i$} ہو تو تصدیق کیجیے کہ:}
\mline{\overline{\left(\dfrac{z_{1}}{z_{2}}\right)}=\dfrac{\overline{z_{1}}}{\overline{z_{2}}}}
\uq{13}{اگر \lr{$z=4-3i$} ہو تو تصدیق کیجیے کہ:}
\mline{|z|=|-z|=\left|\overline{\overline{z}}\right|=|-\overline{z}|}
\uq{14}{حل کریں:}
\mline{\overline{z_{1}z_{2}}=\overline{z_{1}}\ \overline{z_{2}}}
\uq{15}{حل کریں:}
\mline{\overline{\left(\dfrac{z_{1}}{z_{2}}\right)}=\dfrac{\overline{z_{1}}}{\overline{z_{2}}}\ ,\ (z_{2}\neq0)}
\uq{16}{اگر \lr{$z_{1}=5+4i,\ z_{2}=3+2i$} ہوں تو ثابت کیجیے کہ}
\mline{\overline{\left(\dfrac{z_{1}}{z_{2}}\right)}=\dfrac{\overline{z_{1}}}{\overline{z_{2}}}}
\uq{17}{حل کریں:}
\mline{|z|=|-z|=\left|\overline{z}\right|=|-\overline{z}|}
\uq{18}{اگر \lr{$z_{1}=3+4i,\ z_{2}=2+3i$} ہوں تو ثابت کیجیے کہ}
\mline{\overline{z_{1}z_{2}}=\overline{z_{1}}\ \overline{z_{2}}}
\uq{19}{کمپلیکس عددی سروں \lr{z} اور \lr{w} کے لیے یک درجی ہمزدہ مساواتوں کو حل کیجیے:}
\mline{3z+(2+i)w=11-i\ ,\ (2-i)z-w=-1+i}
\uq{20}{کمپلیکس عددی سروں \lr{z} اور \lr{w} کے لیے یک درجی ہمزدہ مساواتوں کو حل کیجیے:}
\mline{2z+(3+i)w=9-i\ ,\ -iz-iw=-1+i}
\uq{21}{کمپلیکس عددی سروں \lr{z} اور \lr{w} کے لیے یک درجی ہمزدہ مساواتوں کو حل کیجیے:}
\mline{z-4w=3i\ ,\ 2z+3w=11-5i}
\uq{22}{درج ذیل کمپلیکس عدد کے حقیقی اور خیالاتی حصہ معلوم کیجیے:}
\mline{\left(\dfrac{2-i}{2+i}\right)^{-2}}
\uq{23}{درج ذیل کمپلیکس عدد کے حقیقی اور خیالاتی حصہ معلوم کیجیے:}
\mline{\left(\dfrac{1-2i}{1+i}\right)^{2}}
\uq{24}{کمپلیکس عددی سروں \lr{z} اور \lr{w} کے لیے یک درجی ہمزدہ مساواتوں کو حل کیجیے:}
\mline{z+w=3i\ ,\ 2z+3w=2}
\uq{25}{کمپلیکس عددی سروں \lr{z} اور \lr{w} کے لیے یک درجی ہمزدہ مساواتوں کو حل کیجیے:}
\mline{2z+(3+i)w=1\ ,\ -z-(1-i)w=2}
\uq{26}{اگر \lr{$z_{1}=3+4i$} اور \lr{$z_{2}=2+3i$} ہو تو تصدیق کیجیے کہ:}
\mline{\overline{\left(\dfrac{z_{1}}{z_{2}}\right)}=\dfrac{\overline{z_{1}}}{\overline{z_{2}}}}
\uq{27}{کمپلیکس سروں \lr{z} اور \lr{w} کے لیے دی گئی یک درجی ہمزاد مساواتوں کو حل کیجیے:}
\mline{iz+(2-i)w=4+i\ ,\ iz+(3+i)w=3+3i}
\uq{28}{\lr{$(3-4i)(a+bi)=1+0i$} کو حل کیجیے اور \lr{a} اور \lr{b} کی قیمتیں معلوم کیجیے۔}
\uq{29}{\lr{x} اور \lr{y} کے لیے مساوات کو حل کیجیے۔ \lr{$(3-2i)(x+yi)=2(x-2yi)+2i-1$}}
\uq{30}{اگر \lr{$z_{1}=3+4i$} اور \lr{$z_{2}=2+3i$} ہو تو تصدیق کیجیے کہ:}
\mline{\overline{z_{1}+z_{2}}=\overline{z_{1}}+\overline{z_{2}}}
\uq{31}{اگر \lr{$z_{1}=3+4i$} اور \lr{$z_{2}=2+3i$} ہو تو تصدیق کیجیے کہ:}
\mline{\overline{z_{1}z_{2}}=\overline{z_{1}}\ \overline{z_{2}}}
\uq{32}{اگر \lr{$z_{1}=3+4i$} اور \lr{$z_{2}=2+3i$} ہو تو تصدیق کیجیے کہ:}
\mline{|z_{1}|=\left|-\overline{z_{1}}\right|}
\uq{33}{اگر \lr{$z_{1}=3+4i$} اور \lr{$z_{2}=2+3i$} ہو تو تصدیق کیجیے کہ:}
\mline{\overline{\overline{z_{2}}}=z_{2}}
\uq{34}{اگر \lr{$z_{1}=3+4i$} اور \lr{$z_{2}=2+3i$} ہو تو تصدیق کیجیے کہ:}
\mline{z_{1}\overline{z_{1}}=|z_{1}|^{2}}
\uq{35}{درج ذیل دو درجی مساوات کو تجزی کے طریقے سے حل کیجیے: \lr{$x^{2}-\dfrac{3}{2}x=\dfrac{9}{2}$}}
\uq{36}{درج ذیل دو درجی مساوات کو تجزی کے طریقے سے حل کیجیے: \lr{$\dfrac{3x-8}{x-2}=\dfrac{5x-2}{x+5},x\neq2,-5$}}
\uq{37}{درج ذیل دو درجی مساوات کو تجزی کے طریقے سے حل کیجیے:}
\mline{\dfrac{1}{x-1}-\dfrac{1}{x+3}=\dfrac{1}{35},x\neq1,-3}
\uq{38}{درج ذیل دو درجی مساوات کو تکمیل مربع کے طریقہ سے حل کیجیے: \lr{$2x^{2}+5x+2=0$}}
\uq{39}{درج ذیل دو درجی مساوات کو تکمیل مربع کے طریقہ سے حل کیجیے: \lr{$x^{2}+x=42$}}
\uq{40}{درج ذیل دو درجی مساوات کو تکمیل مربع کے طریقہ سے حل کیجیے: \lr{$12x^{2}+7x=12$}}
\uq{41}{درج ذیل دو درجی مساوات کو تکمیل مربع کے طریقہ سے حل کیجیے:}
\mline{\dfrac{x+3}{2x-7}=\dfrac{2x-1}{x-3},x\neq\dfrac{7}{2},3}
\uq{42}{درج ذیل دو درجی مساوات کو تکمیل مربع کے طریقہ سے حل کیجیے:}
\mline{\dfrac{1}{1+x}-\dfrac{1}{3-x}=\dfrac{6}{35},x\neq-1,3}
\uq{43}{درج ذیل دو درجی مساوات کو تکمیل مربع کے طریقہ سے حل کیجیے:}
\mline{\dfrac{3x-1}{4x+7}=1-\dfrac{6}{x+7},x\neq\dfrac{-7}{4},-7}
\uq{44}{درج ذیل مساوات کو دو درجی کلیہ کی مدد سے حل کیجیے: \lr{$\dfrac{x+4}{x-4}+\dfrac{x-2}{x-3}=6\dfrac{1}{3},x\neq4,3$}}
\uq{45}{درج ذیل مساوات کو دو درجی کلیہ کی مدد سے حل کیجیے: \lr{$\dfrac{3x-3}{x+1}=\dfrac{2x-1}{x-1},x\neq-1,1$}}
\uq{46}{درج ذیل دو درجی مساوات کو گراف کے طریقے سے حل کیجیے۔ \lr{$x^{2}-3x-18=0$}}
\uq{47}{درج ذیل دو درجی مساوات کو گراف کے طریقے سے حل کیجیے۔ \lr{$x^{2}-5x-14=0$}}
\uq{48}{درج ذیل دو درجی مساوات کو گراف کے طریقے سے حل کیجیے۔ \lr{$2x^{2}+13x+6=0$}}
\uq{49}{درج ذیل دو درجی مساوات کو گراف کے طریقے سے حل کیجیے۔ \lr{$4x^{2}+12x-27=0$}}
\uq{50}{تجزی کے طریقے سے حل کیجیے۔ \lr{$\dfrac{2x-3}{2}=\dfrac{4x-6}{x}$} جبکہ \lr{x} صفر کے برابر نہیں ہے۔}
\uq{51}{تکمیل مربع کے طریقے سے حل کیجیے: \lr{$3x^{2}-32=-10x$}}
\uq{52}{دو درجی تفاعل \lr{$y=2x^{2}-x-10$} کا گراف بنائیے اور حل کیجیے۔}
\uq{53}{درج ذیل یک درجی مساوات کے جوڑوں کو بذریعہ گراف حل کیجیے: \lr{$x+y=8\ ,\ 3x-y=4$}}
\uq{54}{درج ذیل یک درجی مساوات کے جوڑوں کو بذریعہ گراف حل کیجیے: \lr{$x-y=1\ ,\ x+2y=7$}}
\uq{55}{درج ذیل یک درجی مساوات کے جوڑوں کو بذریعہ گراف حل کیجیے:}
\mline{x-2y=1\ ,\ 2x+y=2}
\uq{56}{درج ذیل یک درجی مساوات کے جوڑوں کو بذریعہ گراف حل کیجیے:}
\mline{y=2x+2\ ,\ 3x+2y=4}
\uq{57}{درج ذیل یک درجی مساوات کے جوڑوں کو بذریعہ گراف حل کیجیے:}
\mline{3y=2x+8\ ,\ x+y=1}
\uq{58}{مساوات کے جوڑوں کو بذریعہ گراف حل کیجیے: \lr{$y=8x-32\ ,\ y=x^{2}-6x+8$}}
\uq{59}{مساوات کے جوڑوں کو بذریعہ گراف حل کیجیے: \lr{$y+x=2\ ,\ y=2x^{2}+x-10$}}
\uq{60}{اگر \lr{$x+y=6$} اور \lr{$3x+y=12$} ہوں تو ان کا نقطہ تقاطع بذریعہ گراف معلوم کیجیے۔}
\uq{61}{اگر \lr{$y=2+3x$} اور \lr{$y=2x^{2}+7x-4$} ہوں تو تقاطع کے نقاط بذریعہ گراف معلوم کیجیے۔}
\uq{62}{دو درجی مساوات معلوم کیجیے جس کے روٹس مساوات \lr{$x^{2}+9x+20=0$} کے روٹس سے \lr{2} زیادہ ہوں۔}
\uq{63}{وہ مساوات معلوم کیجیے جس کی اصل مساوات \lr{$x^{2}-px+q=0$} کے روٹس سے دوگنا ہوں۔}
\uq{64}{اگر \lr{$\alpha,\beta$} مساوات \lr{$x^{2}+2x+4=0$} کے روٹس ہوں تو، وہ مساوات معلوم کیجیے۔ جس کے روٹس درج ذیل ہوں:}
\mline{2\alpha-\dfrac{1}{\beta},\ 2\beta-\dfrac{1}{\alpha}}
\uq{65}{اگر \lr{$\alpha,\beta$} مساوات \lr{$x^{2}+2x+4=0$} کے روٹس ہوں تو، وہ مساوات معلوم کیجیے۔ جس کے روٹس درج ذیل ہوں:}
\mline{2\alpha-1,\ 2\beta-1}
\uq{66}{اگر \lr{$\alpha,\beta$} مساوات \lr{$x^{2}+ax+b=0$} کے روٹس ہوں اور \lr{$\alpha^{2},\beta^{2}$} مساوات \lr{$x^{2}+Ax+B=0$} کے روٹس ہوں تو ثابت کیجیے کہ \lr{$A=2b-a^{2},\ B=b^{2}$}۔}
\uq{67}{اگر \lr{$\alpha,\beta$} مساوات \lr{$x^{2}+px+q=0$} کے روٹس ہوں تو شرط معلوم کیجیے کہ: \lr{$\alpha=\beta$}}
\uq{68}{اگر \lr{$\alpha,\beta$} مساوات \lr{$x^{2}+px+q=0$} کے روٹس ہوں تو شرط معلوم کیجیے کہ: \lr{$\alpha=\dfrac{1}{\beta}$}}
\uq{69}{\lr{k} کی قیمت معلوم کیجیے جب کہ مساوات \lr{$x^{2}-(2k+4)x+(7k+1)=0$} ایک ایک روٹ \lr{3} ہو۔}
\uq{70}{\lr{m} کی قیمت معلوم کیجیے۔ کب کہ \lr{$2x^{2}+3x+m=0$} کے روٹس کا مجموعہ روٹس کے حاصل ضرب کا دوگنا ہو۔}
\uq{71}{اگر \lr{$\alpha,\beta$} مساوات \lr{$x^{2}+2x+4=0$} کے روٹس ہوں تو، وہ مساوات معلوم کیجیے۔ جس کے روٹس درج ذیل ہوں:}
\mline{\alpha^{2},\beta^{2}}
\uq{72}{شرط معلوم کیجیے کے مساوات \lr{$ax^{2}+bx+c=0$} کے روٹس کے ایک دوسرے کے ضربی معکوس ہوں۔}
\uq{73}{اگر \lr{$\alpha,\beta$} مساوات \lr{$x^{2}-7x+10=0$} کے روٹس ہوں تو ایک ایسی مساوات بنائیے جس کے روٹس \lr{$2\alpha+1$} اور \lr{$2\beta+1$} ہوں۔}
\uq{74}{اگر دو درجی مساوات \lr{$16x^{2}+7px+49=0$} کے روٹس مساوی ہوں تو \lr{p} کی قیمت معلوم کریں۔}
\uq{75}{اگر دو درجی مساوات \lr{$4u^{2}+8u+q=0$} کے روٹس غیر مساوی اور حقیقی ہوں تو \lr{q} کی ممکنہ قیمتیں معلوم کریں۔}
\uq{76}{\lr{m} کی قیمت معلوم کیجیے اگر دو درجی مساوات \lr{$mx^{2}-8x+1=0$} کے روٹس حقیقی اور مساوی ہوں۔}
\uq{77}{\lr{t} کی کن قیمتوں کے لیے مساوات \lr{$3x^{2}+x+9t=0$} کے روٹس حقیقی اور غیر مساوی ہوں گے۔}
\uq{78}{اگر مساوات \lr{$(k+1)x^{2}+(2k-1)x+(k-1)=0$} کے روٹس برابر ہوں تو \lr{k} کی قیمت معلوم کیجیے۔}
\uq{79}{درج ذیل غیر مساوات کو حل کیجیے: \lr{$x^{2}+3x-4>0$}}
\uq{80}{درج ذیل غیر مساوات کو حل کیجیے: \lr{$2x^{2}-8x+6>0$}}
\uq{81}{درج ذیل غیر مساوات کو حل کیجیے: \lr{$x^{2}+x-6<0$}}
\uq{82}{درج ذیل غیر مساوات کو حل کیجیے: \lr{$x^{2}-6x+9<0$}}
\uq{83}{درج ذیل غیر مساوات کو حل کیجیے: \lr{$4x^{2}-16x+15\leq0$}}
\uq{84}{درج ذیل غیر مساوات کو حل کیجیے: \lr{$-x^{2}+3x-2\geq0$}}
\uq{85}{درج ذیل غیر مساوات کو حل کریں: \lr{$x^{2}-5x+6>0$}}
\uq{86}{درج ذیل غیر مساوات کو حل کریں: \lr{$2x^{2}-5x+2\leq0$}}
\uq{87}{ایک شہر کی آبادی کو مساوات \lr{$P(t)=-2t^{2}+40t+800$} سے ماڈل کیا گیا ہے، جہاں \lr{t} بمطابق \lr{2020} سالوں میں ہے۔ وہ سال معلوم کریں جب آبادی کم از کم \lr{1000} ہوگی۔}
\uq{88}{ایک کھلونا کار ڈھلوان پر پہنچ جاتی ہے جب کہ اس کا طے کردہ فاصلے \lr{(d)} کو مساوات \lr{$d=t^{2}-0.5t$} سے ظاہر کیا جاتا ہے، جہاں وقت \lr{t} سکینڈ میں ہے، وہ وقت معلوم کریں جب کار نے \lr{12.5} میٹر کا فاصلہ طے کیا ہو۔}
\uq{89}{ایک گیند کی بلندی (میٹروں میں) \lr{t} سیکنڈوں کے بعد \lr{$h(t)=-4t^{2}+24t$} ہے۔ کس وقت زمین سے کم از کم \lr{20} میٹر اوپر ہوگی؟}
\uq{90}{ایک فری لانسر کی کمائی \lr{E} (روپے میں) فی ہفتہ \lr{h} (گھنٹوں میں) کو ماڈل \lr{$E(h)=-2h^{2}+40h$} سے ظاہر کیا جاتا ہے، سب سے زیادہ کمائی کے لیے اسے زیادہ سے زیادہ گھنٹے کام کرنا چاہیے۔}
\uq{91}{ایک کمپنی اپنے منافع \lr{(P)} (روپے ہزاروں میں) کو دو درجی مساوات \lr{$P(x)=-2x^{2}+80x-500$} سے ماڈل کرتی ہے۔ جہاں \lr{x} کس چیز کی قیمت روپوں میں ہے۔ وہ قیمت معلوم کریں جو سب سے زیادہ منافع دیتی ہے اور سب سے زیادہ منافع بھی معلوم کریں۔}
\uq{92}{ایک گیند کو \lr{$20ms^{-1}$} کی بلندی ولاسٹی سے اوپر کی طرف پھینکا جاتا ہے۔ زمین پر واپس آنے کا وقت اور زمین سے زیادہ سے زیادہ بلندی معلوم کیجیے۔}
\uq{93}{ایک گیند کو زمین سے اوپر کی طرف پھینکا جاتا ہے۔ \lr{t} سیکنڈ کے بعد اس کی بلندی (میٹر میں) ماڈل کی گئی \lr{$h(t)=-5t^{2}+20t$} ہے۔ کس وقت زمین سے کم از کم \lr{15} میٹر اوپر ہوگی؟}
\uq{94}{درج ذیل دو درجی مساوات کو بذریعہ تجزی، تکمیل مربع اور دو درجی کلیہ حل کیجیے: \lr{$8x^{2}=x+7$}}
\uq{95}{درج ذیل دو درجی مساوات کو بذریعہ تجزی، تکمیل مربع اور دو درجی کلیہ حل کیجیے:}
\mline{2x^{2}-x-10=0}
\uq{96}{اگر ایک گیند کو \lr{v} ولاسٹی سے اوپر پھینکا جائے تو زیادہ سے زیادہ بلندی \lr{(h)} ہے جس پر وہ پہنچ سکتی ہے اور اسے کلیہ \lr{$h=\dfrac{v^{2}}{2g}$} کے ذریعت معلوم کی جاسکتی ہے۔ کلیہ کو اس طرح ترتیب دیں کہ \lr{v} کو سبجیکٹ بنایا جاسکے۔}
\uq{97}{اگر مساوات \lr{$x^{2}+2(1+k)x+k^{2}=0$} کے روٹس مساوی ہوں تو \lr{k} کی قیمت معلوم کریں۔}
\uq{98}{اگر \lr{$\begin{bmatrix}a+2&c-3\\b-1&d+4\end{bmatrix}=\begin{bmatrix}5&8\\6&4\end{bmatrix}$} ہو تو، \lr{a,b,c} اور \lr{d} کی قیمتیں معلوم کیجیے۔}
\uq{99}{اگر \lr{$\begin{bmatrix}a+b&2d-1\\3b+2&4\end{bmatrix}=\begin{bmatrix}10&5\\11&c\end{bmatrix}$} ہو تو \lr{a,b,c} اور \lr{d} کی قیمتیں معلوم کیجیے۔}
\uq{100}{اگر \lr{$\begin{bmatrix}p+q&5\\11&p-2q\end{bmatrix}=\begin{bmatrix}6&5\\11&0\end{bmatrix}$} ہو تو \lr{p} اور \lr{q} کی قیمتیں معلوم کیجیے۔}
\uq{101}{ثابت کیجیے کہ \lr{$L=\begin{bmatrix}2&3&4\\3&-2&5\\4&5&0\end{bmatrix}$} ایک سمیٹرک قالب ہے۔}
\uq{102}{اگر \lr{$A=\begin{bmatrix}2&3\\-3&2\end{bmatrix},\ B=\begin{bmatrix}3&4\\5&6\end{bmatrix}$} اور \lr{$C=\begin{bmatrix}1&-2\\0&5\end{bmatrix}$} ہو تو درج ذیل کی تصدیق کریں:}
\mline{A+B=B+A}
\uq{103}{اگر \lr{$A=\begin{bmatrix}2&3\\-3&2\end{bmatrix},\ B=\begin{bmatrix}3&4\\5&6\end{bmatrix}$} اور \lr{$C=\begin{bmatrix}1&-2\\0&5\end{bmatrix}$} ہو تو درج ذیل کی تصدیق کریں:}
\mline{(A+B)+C=A+(B+C)}
\uq{104}{اگر \lr{$A=\begin{bmatrix}2&3\\-3&2\end{bmatrix},\ B=\begin{bmatrix}3&4\\5&6\end{bmatrix}$} اور \lr{$C=\begin{bmatrix}1&-2\\0&5\end{bmatrix}$} ہو تو درج ذیل کی تصدیق کریں:}
\mline{(2A+B)+C=2A+(B+C)}
\uq{105}{اگر \lr{$A=\begin{bmatrix}2&3\\-3&2\end{bmatrix},\ B=\begin{bmatrix}3&4\\5&6\end{bmatrix}$} اور \lr{$C=\begin{bmatrix}1&-2\\0&5\end{bmatrix}$} ہو تو درج ذیل کی تصدیق کریں:}
\mline{3(A+B)=3A+3B}
\uq{106}{اگر \lr{$A=\begin{bmatrix}5&6\\7&-2\end{bmatrix}$} اور \lr{$B=\begin{bmatrix}-5&-6\\-7&2\end{bmatrix}$} ہو تو ثابت کیجیے کہ قالب \lr{A} قالب \lr{B} کا اور \lr{B} قالب \lr{A} کا جمعی معکوس ہے۔}
\uq{107}{اگر \lr{$A=\begin{bmatrix}6&-2\\0&3\end{bmatrix}$} اور \lr{$B=\begin{bmatrix}8&-1\\3&0\end{bmatrix}$} ہوں تو تصدیق کیجیے: \lr{$(A+B)^{t}=A^{t}+B^{t}$}}
\uq{108}{اگر \lr{$A=\begin{bmatrix}6&-2\\0&3\end{bmatrix}$} اور \lr{$B=\begin{bmatrix}8&-1\\3&0\end{bmatrix}$} ہوں تو تصدیق کیجیے: \lr{$(A-B)^{t}=A^{t}-B^{t}$}}
\uq{109}{اگر \lr{$A=\begin{bmatrix}3&8&2\\5&3&1\\6&0&7\end{bmatrix},\ B=\begin{bmatrix}4&1&-2\\3&0&6\\2&3&4\end{bmatrix}$} ہوں تو \lr{$A+B$} اور \lr{$B+A$} معلوم کیجیے۔}
\uq{110}{اگر \lr{$A=\begin{bmatrix}5&9\\4&8\end{bmatrix},\ B=\begin{bmatrix}3&6\\2&1\end{bmatrix}$} ہوں تو \lr{$A-B$} اور \lr{$B-A$} معلوم کیجیے۔}
\uq{111}{\lr{$A=\begin{bmatrix}5&1\\-1&4\end{bmatrix},\ B=\begin{bmatrix}2&-1\\0&3\end{bmatrix}$} اور \lr{$C=\begin{bmatrix}5&1\\1&4\end{bmatrix}$} ہوں تو ثابت کیجیے: \lr{$AB\neq BA$}}
\uq{112}{\lr{$A=\begin{bmatrix}5&1\\-1&4\end{bmatrix},\ B=\begin{bmatrix}2&-1\\0&3\end{bmatrix}$} اور \lr{$C=\begin{bmatrix}5&1\\1&4\end{bmatrix}$} ہوں تو ثابت کیجیے:}
\mline{A(B-C)=AB-AC}
\uq{113}{\lr{$A=\begin{bmatrix}5&1\\-1&4\end{bmatrix},\ B=\begin{bmatrix}2&-1\\0&3\end{bmatrix}$} اور \lr{$C=\begin{bmatrix}5&1\\1&4\end{bmatrix}$} ہوں تو ثابت کیجیے:}
\mline{A(BC)=(AB)C}
\uq{114}{\lr{$A=\begin{bmatrix}5&1\\-1&4\end{bmatrix},\ B=\begin{bmatrix}2&-1\\0&3\end{bmatrix}$} اور \lr{$C=\begin{bmatrix}5&1\\1&4\end{bmatrix}$} ہوں تو ثابت کیجیے:}
\mline{(BC)^{t}=C^{t}B^{t}}
\uq{115}{\lr{$A=\begin{bmatrix}5&1\\-1&4\end{bmatrix},\ B=\begin{bmatrix}2&-1\\0&3\end{bmatrix}$} اور \lr{$C=\begin{bmatrix}5&1\\1&4\end{bmatrix}$} ہوں تو ثابت کیجیے:}
\mline{(B+C)A=BA+CA}
\uq{116}{اگر \lr{$A=\begin{bmatrix}5&-3\\2&-1\end{bmatrix}$} ہو تو \lr{$A^{-1}$} معلوم کیجیے اور ثابت کیجیے کہ \lr{$AA^{-1}=A^{-1}A=I$}}
\uq{117}{ثابت کیجیے کہ درج ذیل قالب ایک دوسرے کے ضربی معکوس ہیں: \lr{$\begin{bmatrix}2&-1\\-3&2\end{bmatrix},\begin{bmatrix}2&1\\3&2\end{bmatrix}$}}
\uq{118}{ثابت کیجیے کہ درج ذیل قالب ایک دوسرے کے ضربی معکوس ہیں: \lr{$\begin{bmatrix}4&15\\2&8\end{bmatrix},\begin{bmatrix}4&-\frac{15}{2}\\-1&2\end{bmatrix}$}}
\uq{119}{ثابت کیجیے کہ \lr{$(AB)^{-1}=B^{-1}A^{-1}$} اگر \lr{$A=\begin{bmatrix}-3&-2\\5&6\end{bmatrix},\ B=\begin{bmatrix}2&-1\\-3&2\end{bmatrix}$}}
\uq{120}{ثابت کیجیے کہ \lr{$(AB)^{-1}=B^{-1}A^{-1}$} اگر \lr{$A=\begin{bmatrix}1&2\\8&0\end{bmatrix},\ B=\begin{bmatrix}0&-1\\5&2\end{bmatrix}$}}
\uq{121}{اگر \lr{$A=\begin{bmatrix}8&5\\4&3\end{bmatrix}$} ہو تو \lr{A} کا ضربی معکوس معلوم کیجیے اور ثابت کیجیے کہ}
\uq{122}{ثابت کیجیے کہ \lr{$A=\begin{bmatrix}4&-2\\-6&4\end{bmatrix}$} اور \lr{$B=\begin{bmatrix}1&\frac{1}{2}\\\frac{3}{2}&1\end{bmatrix}$} ایک دوسرے کے ضربی معکوس ہیں۔}
\uq{123}{اگر ممکن ہو تو معکوس قالب کے طریقے سے حل کیجیے:}
\mline{2x+5y=19\ ,\ 4x-3y=-1}
\uq{124}{اگر ممکن ہو تو معکوس قالب کے طریقے سے حل کیجیے: \lr{$3x+2y=7\ ,\ 5x-y=16$}}
\uq{125}{اگر ممکن ہو تو معکوس قالب کے طریقے سے حل کیجیے: \lr{$x-2y=9\ ,\ 2x+7y=-4$}}
\uq{126}{اگر ممکن ہو تو معکوس قالب کے طریقے سے حل کیجیے: \lr{$3x+2y=2\ ,\ x-2y=-2$}}
\uq{127}{اگر ممکن ہو تو درج ذیل یک درجی مساوات کے جوڑوں کو کریمر کے قانون کی مدد سے حل کریں:}
\mline{x+4y=4\ ,\ 2x-y=5}
\uq{128}{اگر ممکن ہو تو درج ذیل یک درجی مساوات کے جوڑوں کو کریمر کے قانون کی مدد سے حل کریں:}
\mline{x+2y=7\ ,\ 3x-2y=-3}
\uq{129}{اگر ممکن ہو تو درج ذیل یک درجی مساوات کے جوڑوں کو کریمر کے قانون کی مدد سے حل کریں:}
\mline{2x-5y=-6\ ,\ 4x-3y=-12}
\uq{130}{اگر ممکن ہو تو درج ذیل یک درجی مساوات کے جوڑوں کو کریمر کے قانون کی مدد سے حل کریں:}
\mline{3x+2y=-1\ ,\ 5x+6y=5}
\uq{131}{ایک الیکٹریکل انجینئر ایک سادہ برقی سرکٹ کی دو شاخوں \lr{A} اور \lr{B} میں کرنٹ معلوم کرنا چاہتا ہے۔ مساواتوں کا نظام یہ ہے: \lr{$x+y=7\ ,\ 2x-y=2$} جہاں \lr{x} شاخ \lr{A} میں کرنٹ ہے \lr{y} شاخ \lr{B} میں کرنٹ ہے۔ قالبوں کا استعمال کرتے ہوئے \lr{x} اور \lr{y} کی قیمت معلوم کیجیے۔}
\uq{132}{تین قوتیں ایک ذرہ پر عمل کر رہی ہیں جو کہ متوازن ہوں۔ یعنی \lr{$F_{1}+F_{2}+F_{3}=0$}، جہاں \lr{$F_{1}=\begin{bmatrix}8\\x\end{bmatrix},\ F_{2}=\begin{bmatrix}-2\\-7\end{bmatrix},\ F_{3}=\begin{bmatrix}y\\-1\end{bmatrix}$} ہے۔ \lr{x} اور \lr{y} کی قیمت معلوم کریں۔}
\uq{133}{دو سپورٹ بیمز \lr{A} اور \lr{B} \lr{100} کلونیوٹن کا مشترکہ بوجھ اٹھا رہے ہیں، بیم \lr{A} پر بوجھ کا دوگنا اور بیم \lr{B} پر بوجھ کا تین گنا \lr{240} کلونیوٹن بنتا ہے۔ قالبوں کا استعمال کرتے ہوئے بیم \lr{A} اور بیم \lr{B} پر بوجھ معلوم کریں۔}
\uq{134}{ایک \lr{D2} گیم کی دنیا میں، دو کردار خط مستقیم پر چل رہے ہیں۔ ایک کردار ایسے خط پر چلتا ہے جہاں اس کی افقی پوزیشن اور عمودی پوزیشن کے دوگنا کا مجموعہ \lr{5} ہے۔ جبکہ دوسرا کردار ایسے خط پر چلتا ہے جہاں اس کی افقی پوزیشن اس کی عمودی پوزیشن سے ایک زیادہ ہے۔ قالبوں کا استعمال کرتے ہوئے ان کا نقطہ تقاطع معلوم کیجیے۔}
\uq{135}{دو سال پہلے ایک آدمی اپنے بیٹے سے \lr{5} گنا بڑا تھا۔ \lr{6} سال بعد وہ اپنے بیٹے سے \lr{3} گنا بڑا ہو گیا۔ قالبوں کا استعمال کرتے ہوئے ان کی موجودہ عمریں معلوم کیجیے۔}
\uq{136}{دو سائیکل سوار \lr{44} کلو میٹر کے فاصلے پر ہیں اور ایک ہی وقت میں سفر شروع کرتے ہیں۔ اگر وہ ایک دوسرے کی طرف چلیں تو وہ دو گھنٹے میں ملتے ہیں لیکن اگر وہ ایک ہی سمت میں چلیں تو تیز رفتار سائیکل سوار سست رفتار سائیکل سوار کو \lr{$7\frac{1}{2}$} گھنٹے میں پیچھے چھوڑ دیتا ہے۔ قالبوں کا استعمال کرتے ہوئے ان کی رفتاریں معلوم کریں۔}
\uq{137}{ایک کسر کا شمار کنندہ اس کے مخرج سے \lr{7} کم ہے۔ اگر شمار کنندہ میں \lr{3} کا اضافہ کر دیا جائے تو نئی کسر \lr{$\frac{3}{4}$} بن جاتی ہے۔ قالبوں کا استعمال کرتے ہوئے اصل کسر معلوم کیجیے۔}
\uq{138}{درج ذیل کو معکوس قالب کے طریقے سے حل کیجیے: \lr{$2x+3y=13\ ,\ 4x-5y=-7$}}
\uq{139}{کریمر کے قانون کا استعمال کرتے ہوئے مساواتوں کا حل کیجیے:}
\mline{2x+3y=13\ ,\ 4x-5y=-7}
\uq{140}{آٹھ سال پہلے جویریہ کی عمر جنت کی عمر کا \lr{$\frac{3}{4}$} گنا تھی۔ چار سال بعد جویریہ کی عمر جنت کی عمر کا \lr{$\frac{6}{7}$} گنا ہو گی۔ قالبوں کا استعمال کرتے ہوئے ان کی موجودہ عمریں معلوم کیجیے۔}
\uq{141}{تین قوتیں ایک ذرہ پر عمل کر رہی ہیں جو کے لازمی طور پر متوازن ہوں، یعنی \lr{$F_{1}+F_{2}+F_{3}=0$} جہاں \lr{$F_{1}=\begin{bmatrix}5\\2\end{bmatrix},\ F_{2}=\begin{bmatrix}-3\\x\end{bmatrix},\ F_{3}=\begin{bmatrix}-2\\-4\end{bmatrix}$} ہے، \lr{x} کی قیمت معلوم کیجیے۔}
\uq{142}{اگر \lr{$A=\begin{bmatrix}2&3\\4&5\end{bmatrix},\ B=\begin{bmatrix}5&6\\2&1\end{bmatrix},\ C=\begin{bmatrix}3&2\\1&7\end{bmatrix}$} ہوں تو تصدیق کیجیے:}
\mline{2(A+B)=2A+2B}
\uq{143}{اگر \lr{$A=\begin{bmatrix}2&3\\4&5\end{bmatrix},\ B=\begin{bmatrix}5&6\\2&1\end{bmatrix},\ C=\begin{bmatrix}3&2\\1&7\end{bmatrix}$} ہوں تو تصدیق کیجیے:}
\mline{(A+B)+C=A+(B+C)}
\uq{144}{اگر \lr{$A=\begin{bmatrix}2&3\\4&5\end{bmatrix},\ B=\begin{bmatrix}5&6\\2&1\end{bmatrix},\ C=\begin{bmatrix}3&2\\1&7\end{bmatrix}$} ہوں تو تصدیق کیجیے:}
\mline{(A+B)C=AC+BC}
\uq{145}{اگر \lr{$A=\begin{bmatrix}2&3\\4&5\end{bmatrix},\ B=\begin{bmatrix}5&6\\2&1\end{bmatrix},\ C=\begin{bmatrix}3&2\\1&7\end{bmatrix}$} ہوں تو تصدیق کیجیے:}
\mline{C(A-B)=CA-CB}
\uq{146}{اگر \lr{$A=\begin{bmatrix}2&3\\4&5\end{bmatrix},\ B=\begin{bmatrix}5&6\\2&1\end{bmatrix},\ C=\begin{bmatrix}3&2\\1&7\end{bmatrix}$} ہوں تو تصدیق کیجیے:}
\mline{(AB)^{-1}=B^{-1}A^{-1}}
\uq{147}{اگر \lr{$A=\begin{bmatrix}2&3\\4&5\end{bmatrix},\ B=\begin{bmatrix}5&6\\2&1\end{bmatrix},\ C=\begin{bmatrix}3&2\\1&7\end{bmatrix}$} ہوں تو تصدیق کیجیے:}
\mline{AA^{-1}=A^{-1}A=1}
\uq{148}{اگر \lr{$A=\begin{bmatrix}2&3\\4&5\end{bmatrix},\ B=\begin{bmatrix}5&6\\2&1\end{bmatrix},\ C=\begin{bmatrix}3&2\\1&7\end{bmatrix}$} ہوں تو تصدیق کیجیے:}
\mline{(AB)^{t}=B^{t}A^{t}}
\uq{149}{اگر \lr{$A=\begin{bmatrix}2&3\\4&5\end{bmatrix},\ B=\begin{bmatrix}5&6\\2&1\end{bmatrix},\ C=\begin{bmatrix}3&2\\1&7\end{bmatrix}$} ہوں تو تصدیق کیجیے:}
\mline{(AB)C=A(BC)}
\uq{150}{اگر ممکن ہو تو درج ذیل یک درجی مساوات کے جوڑے کو معکوس قالب کے طریقے اور کریمر کی مدد سے حل کیجیے:}
\mline{3x+4y=7\ ,\ 5x-y=2}
\uq{151}{اگر ممکن ہو تو درج ذیل یک درجی مساوات کے جوڑے کو معکوس قالب کے طریقے اور کریمر کی مدد سے حل کیجیے:}
\mline{x-6y=-15\ ,\ 2x+6y=-3}
\uq{152}{اگر ممکن ہو تو درج ذیل یک درجی مساوات کے جوڑے کو معکوس قالب کے طریقے اور کریمر کی مدد سے حل کیجیے:}
\mline{2x+y=5\ ,\ x+3y=3}
\uq{153}{قالبوں کا استعمال کرتے ہوئے دو ایسے اعداد معلوم کیجیے کہ پہلے عدد کا دوگنا دوسرے میں شامل کرنے سے \lr{21} بنتا ہے اور دوسرے عدد کا دوگنا پہلے میں جمع کرنے سے \lr{27} بنتا ہے۔}
\uq{154}{\lr{4} چاقو اور \lr{6} کانٹوں کی قیمت \lr{136} روپے ہے، جب کہ \lr{6} چاقوں اور \lr{5} کانٹوں کی قیمت \lr{164} روپے ہے۔ قالبوں کا استعمال کرتے ہوئے ایک چاقو اور ایک کانٹے کی قیمت معلوم کیجیے۔}
\uq{155}{ایک دوکاندار \lr{5} مردوں اور \lr{3} عورتوں کو ملازمت دیتا ہے۔ اور \lr{3500} روپے یومیہ اجرت ادا کرتا ہے۔ اگر مردوں کی تعداد کم کرکے \lr{2} کر دی جائے اور \lr{3} اضافی عورتوں کو ملازم رکھا جائے، تو یومیہ اجرت کی رقم \lr{5000} روپے ہو جاتی ہے۔ قالبوں کا استعمال کرتے ہوئے ایک مرد اور ایک عورت کی یومیہ اجرت معلوم کیجیے۔}
\uq{156}{\lr{k} کی قیمت معلوم کیجیے بشرطیکہ \lr{$(fog)(x)=(gof)(x)$}، جہاں \lr{$f(x)=3x+2,g(x)=6x-k$}۔}
\uq{157}{تصدیق کیجیے: \lr{$f\left(f^{-1}(x)\right)=f^{-1}(f(x))=x$} معلوم کیے بغیر: \lr{$f(x)=x-6$}}
\uq{158}{تصدیق کیجیے: \lr{$f\left(f^{-1}(x)\right)=f^{-1}(f(x))=x$} معلوم کیے بغیر: \lr{$f(x)=7x-4$}}
\uq{159}{تصدیق کیجیے: \lr{$f\left(f^{-1}(x)\right)=f^{-1}(f(x))=x$} معلوم کیے بغیر: \lr{$f(x)=\dfrac{x-3}{4}$}}
\uq{160}{تصدیق کیجیے: \lr{$f\left(f^{-1}(x)\right)=f^{-1}(f(x))=x$} معلوم کیے بغیر:}
\mline{f(x)=\dfrac{x-4}{x+2},x\neq-2}
\uq{161}{اگر \lr{$f(x)=x^{2}+9$} اور \lr{$g(x)=x+21$} ہو تو \lr{a} کی قیمتیں معلوم کیجیے جب کہ}
\mline{f(a)=g(a)}
\uq{162}{اگر \lr{$f(x)=x^{3}$} ہو تو \lr{$f^{-1}(x)$} معلوم کریں۔ اور تصدیق کیجیے کہ}
\mline{f\left(f^{-1}(x)\right)=f^{-1}(f(x))=x}
\uq{163}{درج ذیل مطلق قدر والے تفاعل کے گراف بنائیے: \lr{$f(x)=|x-2|$}}
\uq{164}{درج ذیل مطلق قدر والے تفاعل کے گراف بنائیے: \lr{$f(x)=3|x+3|-4$}}
\uq{165}{درج ذیل مطلق قدر والے تفاعل کے گراف بنائیے: \lr{$f(x)=5|x|$}}
\uq{166}{درج ذیل مطلق قدر والے تفاعل کے گراف بنائیے: \lr{$f(x)=|x+2|+3$}}
\uq{167}{درج ذیل مطلق قدر والے تفاعل کے گراف بنائیے: \lr{$f(x)=2|x+4|-3$}}
\uq{168}{درج ذیل مطلق قدر والے تفاعل کے گراف بنائیے: \lr{$f(x)=2|x+1|-6$}}
\uq{169}{حل کیجیے اور حل کو عددی خط پر ظاہر کیجیے: \lr{$1-2\left|\dfrac{3}{2}x-5\right|>-3$}}
\uq{170}{\lr{$f(x)=2|x+1|-3$} کا گراف بنائیے۔}
\uq{171}{ایک تفاعل \lr{$B(t)=5{,}000+200t$} جو \lr{t} مہینوں کے بعد کل بیلنس (روپوں میں) کو ظاہر کرتا ہے۔ \lr{6} مہینوں کے بعد بیلنس کتنا ہوگا؟}
\uq{172}{ایک تفاعل \lr{$f(k)=150+20k$} جو \lr{k} کلو میٹر کی مسافت کے لیے کل کرائے (روپوں میں) کو ظاہر کرتا ہے۔ \lr{12} کلو میٹر کی سواری کے لیے کتنا کرایہ ہوگا؟}
\uq{173}{فینسی صوفہ سیٹ بنانے کے لیے مقررہ لاگت \lr{5500} روپے ہے، جسے \lr{$f(n)=5500\,n$} کے طور پر ماڈل کیا گیا ہے، جہاں دو صوفہ سیٹوں کی تعداد ہے۔ \lr{50} صوفہ سیٹوں کی لاگت معلوم کیجیے۔}
\uq{174}{ایک تفاعل \lr{$T(d)=\dfrac{d}{60}$}، \lr{d} کلو میٹر کا فاصلہ \lr{T} گھنٹوں میں طے کرنے کو ظاہر کرتا ہے۔ \lr{180} کلو میٹر کا فاصلہ طے کرنے میں کتنا وقت لگے گا۔}
\uq{175}{ایک کمپنی انکوڈنگ کے کام کے لیے \lr{100} روپے وصول کرتی ہے اس کے علاوہ کمپنی چھاپے گئے ہر صفحے کے \lr{5} روپے وصول کرتی ہے۔ تفاعل کا ماڈل \lr{$f(x)=100+5x$} ہے۔ جہاں \lr{x} چھاپے گئے صفحات کی تعداد کو ظاہر کرتا ہے۔ کمپنی \lr{55} صفحات کی انکوڈنگ اور پرنٹنگ کے کام کے لیے کتنا معاوضہ وصول کرے گی؟}
\uq{176}{ایک کیمیائی تعامل \lr{$37^{o}C$} پر مستحکم ہو جاتا ہے۔ اس عمل اس وقت ضرور رک جانا چاہیے جب وہ درجہ حرارت میں انحراف \lr{$2.5^{o}C$} سے زیادہ ہو جائے۔ اس حالت کو \lr{$|T-37^{o}|>2.5^{o},\ T$} کے طور پر ماڈل کیا گیا ہے، جہاں \lr{T} درجہ حرارت ہے۔ درجہ حرارت کی کن قیمتوں کے لیے عمل کو روک دینا چاہیے؟}
\uq{177}{ایک فیکٹری دھاتی سلاخیں تیار کرتی ہے جس کی لمبائی \lr{2.5} میٹر ہونی چاہیے جس میں \lr{$\pm0.04$} میٹر کا فرق قابل برداشت ہے۔ اس حالت کو مطلق قدر والی غیر مساوات \lr{$|x-2.5|\leq0.04$} سے ماڈل کیا گیا ہے۔ قابل قبول لمبائیوں کی سعت کیا ہے؟}
\uq{178}{ایک مشین کے پرزے کو اس طرح سیدھ میں کیا جانا چاہیے کہ اس کا مرکز بالکل \lr{0} پر ہو۔ اگر یہ \lr{0.1} ملی میٹر سے زیادہ اپنے مرکز سے ہٹ جائے تو پرزے کو مسترد کر دیا جاتا ہے۔ اس حالت کو \lr{$|x|>0.1$} کے ذریعے ماڈل کیا گیا ہے۔ مرکز کی کون سی حالتیں مسترد سبب ہونے کا سبب بنتی ہیں؟}
\uq{179}{فرض کیا \lr{$E=(n)=1{,}500n$} کل آمدنی (روپوں میں) کا تفاعل ہے، جہاں \lr{n} دونوں کی تعداد ہے۔ اگر ایک مزدور \lr{10} دن کام کرتا ہے تو اس کی آمدنی معلوم کریں۔}
\uq{180}{فرض کیا \lr{$f(x)=200+15x$} ٹیکسی کے کرائے (روپوں میں) کا تفاعل ہے۔ جہاں \lr{x} کلو میٹر کی تعداد کو ظاہر کرتا ہے۔ \lr{50} کلو میٹر کے فاصلے کے لیے مطلوبہ کرایہ معلوم کیجیے۔}
\uq{181}{ایک حساس برقی آلہ \lr{220V} پر کام کرتا ہے۔ لیکن \lr{5V} تک اتار چڑھاؤ کو برداشت کر سکتا ہے۔ یہ حالت \lr{$|V-220|\leq5$} کے ذریعے ماڈل کی گئی ہے قابل قبول وولٹج کی حد کیا ہے؟}
\uq{182}{ایک زلزلہ سینسر اس وقت انتباہ جاری کرتا ہے جب ریکارڈ شدہ شدت محفوظ سطح \lr{4.5} سے \lr{1.2} یونٹس کے فرق کے ساتھ ہو۔ حالت کو \lr{$|M-4.5|>1.2$} کے ذریعے ماڈل کیا گیا ہے۔ شدت \lr{M} کی کن قیمتوں کے لیے سینسر انتباہ جاری کرے گا۔}
\uq{183}{\lr{$f^{-1}(x)$} معلوم کیجیے اگر \lr{$f(x)=\dfrac{5}{x-1},x\neq1$}}
\uq{184}{\lr{$f^{-1}(x)$} معلوم کیجیے اگر \lr{$f(x)=\sqrt{x-5},x\geq5$}}
\uq{185}{\lr{$f^{-1}(x)$} معلوم کیجیے اگر \lr{$f(x)=\dfrac{3-x}{2}$}}
\uq{186}{درج ذیل مطلق قدروں والے تفاعل کا گراف بنائیے: \lr{$f(x)=7|x|$}}
\uq{187}{درج ذیل مطلق قدروں والے تفاعل کا گراف بنائیے: \lr{$f(x)=|x+6|-2$}}
\uq{188}{اگر \lr{$f(x)=3x+7$} ہو تو معلوم کیجیے: \lr{$f^{-1}(x)$}}
\uq{189}{اگر \lr{$f(x)=3x+7$} ہو تو معلوم کیجیے: \lr{x} کی قیمت جس کے لیے \lr{$f(x)=f^{-1}(x)$}}
\uq{190}{ایک کمپنی ہر فروخت شدہ چیز پر \lr{100} روپے منافع کماتی ہے اور اس کا ماہانہ خرچہ \lr{10،000} روپے ہے۔ تفاعل \lr{$P(x)=100x-10{,}000$}، \lr{x} اشیا فروخت کرنے کے منافع کو ظاہر کرتا ہے۔ بریک ایوان (منافع=0) تک پہنچنے کے لیے کتنی اشیا فروخت کرنا ضروری ہے؟}
\uq{191}{ایک سٹور کسی ایک چیز پر \lr{15} فیصد رعایت دیتا ہے۔ تفاعل \lr{$D(p)=0.85p$} اصل قیمت \lr{p} پر رعایت کے بعد کی قیمت فروخت دیتا ہے۔ اس چیز کی قیمت فروخت کیا ہے جس کی اصل لاگت \lr{2،000} روپے ہو۔}
\uq{192}{ایک جی پی ایس سسٹم کو اس وقت درست مانا جاتا ہے۔ اگر اصل مقام رپورٹ کردہ مقام سے \lr{6} میٹر سے زیادہ مختلف نہ ہو۔ اگر اصل مقام نقطہ \lr{x=100} پر ہے۔ تو یہ قابل قبول حد کو اس طرح بیان کیا گیا ہے \lr{$|x-r|\leq6$}، جہاں \lr{r} رپورٹ شدہ مقام ہے۔ قابل قبول رپورٹ کردہ مقامات کی حد کیا ہے؟}
\uq{193}{\lr{$f^{-1}(x)$} معلوم کیجیے اگر \lr{$f(x)=9x-1$}}
\uq{194}{حل کیجیے اور حل کو عددی خط پر ظاہر کیجیے: \lr{$|3x-2|=1$}}
\uq{195}{حل کیجیے اور حل کو عددی خط پر ظاہر کیجیے: \lr{$|6x+1|=9$}}
\uq{196}{حل کیجیے اور حل کو عددی خط پر ظاہر کیجیے: \lr{$|7-2x|\leq1$}}
\uq{197}{مختصر کیجیے: \lr{$\dfrac{4x}{x^{2}-3x+2}-\dfrac{4}{1-x}-\dfrac{5}{x-2}$}}
\uq{198}{مختصر کیجیے: \lr{$\dfrac{x+y}{12x-6y}+\dfrac{x-y}{18x-9y}$}}
\uq{199}{مختصر کیجیے: \lr{$\dfrac{2}{x+2y}-\dfrac{x-6y}{x^{2}-4y^{2}}$}}
\uq{200}{مختصر کیجیے: \lr{$\dfrac{1}{x+1}-\dfrac{2}{x+2}-\dfrac{2x+3}{x^{2}+3x+2}$}}
\uq{201}{مختصر کیجیے: \lr{$\dfrac{5}{x^{2}+x-6}-\dfrac{1}{2x^{2}-7x+6}$}}
\uq{202}{مختصر کیجیے: \lr{$\dfrac{2x}{x+y}-\dfrac{y}{x-y}-\dfrac{2y^{2}}{x^{2}-y^{2}}$}}
\uq{203}{مختصر کیجیے: \lr{$\dfrac{7}{2x^{2}-x-6}-\dfrac{8}{3x^{2}-4x-4}$}}
\uq{204}{مختصر کیجیے: \lr{$\dfrac{x+2}{x^{2}-x-12}-\dfrac{x}{x^{2}+6x+9}$}}
\uq{205}{مختصر کیجیے: \lr{$\dfrac{1}{x^{2}-4x+3}+\dfrac{1}{x^{2}-3x+2}-\dfrac{1}{x^{2}-5x+6}$}}
\uq{206}{\lr{$\dfrac{1}{x^{2}+2}$} کو \lr{$\dfrac{2x^{3}+x^{2}+3}{(x^{2}+2)^{2}}$} میں سے تفریق کیجیے۔}
\uq{207}{مختصر کیجیے: \lr{$\dfrac{x-2y}{4}+\dfrac{x+y}{6}-\dfrac{2x-y}{15}$}}
\uq{208}{مختصر کیجیے: \lr{$\dfrac{4}{x^{2}-2xy-3y^{2}}+\dfrac{1}{x^{2}+3xy+2y^{2}}$}}
\uq{209}{مختصر کیجیے: \lr{$\dfrac{1}{x^{2}-5x+6}+\dfrac{1}{x^{2}-3x+2}-\dfrac{1}{x^{2}-8x+15}$}}
\uq{210}{مختصر کیجیے: \lr{$\dfrac{a^{2}-4b^{2}}{a^{2}+2ba}\times\dfrac{2a^{2}+10ab}{a^{2}+3ab-10b^{2}}$}}
\uq{211}{مختصر کیجیے: \lr{$\dfrac{(a+b)^{2}-c^{2}}{(a+c)^{2}-b^{2}}\times\dfrac{a^{2}-(b-c)^{2}}{(a+c)^{2}-c^{2}}$}}
\uq{212}{مختصر کیجیے: \lr{$\dfrac{x^{2}-x-6}{x^{2}-x-20}\div\left\{\dfrac{x^{3}-3x^{2}}{x^{2}+4x}\times\dfrac{x^{3}-5x^{2}-14x}{x^{2}-12x+35}\right\}$}}
\uq{213}{مختصر کیجیے: \lr{$\dfrac{3x^{2}-3xy}{10x^{2}+10xy-20y^{2}}\times\dfrac{5xy+10y^{2}}{6x^{2}+6xy+6y^{2}}\div\dfrac{2x^{2}-2y^{2}}{4x^{3}-4y^{3}}$}}
\uq{214}{مختصر کیجیے: \lr{$\dfrac{2x^{2}-98}{x^{3}-125}\times\dfrac{x^{2}-3x-10}{3x-21}\div\dfrac{x^{2}+5x-14}{x^{2}+5x+25}$}}
\uq{215}{مختصر کیجیے: \lr{$\dfrac{x^{2}-2x}{2x+6}\times\dfrac{x^{2}+x-6}{x^{2}-5x}\times\dfrac{6x-30}{x^{2}-2x}$}}
\uq{216}{مختصر کیجیے: \lr{$\dfrac{x^{3}-a^{3}}{x^{2}-2bx+b^{2}}\times\dfrac{x^{2}-bx-cx+bc}{ax-x^{2}}\div\dfrac{x^{2}-cx}{x-b}$}}
\uq{217}{درج ذیل مساوات کو حل کیجیے: \lr{$5x^{4}-19x^{2}+12=0$}}
\uq{218}{درج ذیل مساوات کو حل کیجیے: \lr{$4x^{4}-27x^{2}+18=0$}}
\uq{219}{درج ذیل مساوات کو حل کیجیے: \lr{$5x^{4}-22x^{2}+8=0$}}
\uq{220}{درج ذیل مساوات کو حل کیجیے: \lr{$4.2^{2x+1}-9.2^{x}+1=0$}}
\uq{221}{درج ذیل مساوات کو حل کیجیے: \lr{$4^{1+x}+4^{1-x}-10=0$}}
\uq{222}{درج ذیل مساوات کو حل کیجیے: \lr{$3^{x}+3^{3-x}-12=0$}}
\uq{223}{درج ذیل مساوات کو حل کیجیے: \lr{$5^{1+3x}+5^{2-3x}-126=0$}}
\uq{224}{درج ذیل مساوات کو حل کیجیے: \lr{$\left(x^{2}+\dfrac{1}{x^{2}}\right)+2\left(x+\dfrac{1}{x}\right)-33=0$}}
\uq{225}{درج ذیل مساوات کو حل کیجیے: \lr{$2\left(x^{2}+\dfrac{1}{x^{2}}\right)-25\left(x-\dfrac{1}{x}\right)-17=0$}}
\uq{226}{درج ذیل مساوات کو حل کیجیے: \lr{$2x^{4}-5x^{3}-14x^{2}-5x+2=0$}}
\uq{227}{درج ذیل مساوات کو حل کیجیے: \lr{$(x+1)(x+2)(x-4)(x-5)+8=0$}}
\uq{228}{درج ذیل مساوات کو حل کیجیے: \lr{$(x+3)(x+4)(x+5)(x+6)=-1$}}
\uq{229}{حل کیجیے: \lr{$2x^{4}-11x^{2}+12=0$}}
\uq{230}{حل کیجیے: \lr{$4^{x}-3.2^{x+2}+128=0$}}
\uq{231}{حل کیجیے: \lr{$3\left(x^{2}+\dfrac{1}{x^{2}}\right)+4\left(x+\dfrac{1}{x}\right)-14=0$}}
\uq{232}{حل کیجیے: \lr{$(x+2)(x+3)(x-5)(x-6)=-12$}}
\uq{233}{ایک ٹرین یکساں رفتار سے \lr{240} کلو میٹر کا فاصلہ طے کرتی ہے۔ اگر اس نے اپنی رفتار \lr{4} کلو میٹر فی گھنٹہ کم کر دی ہوتی، تو سفر میں \lr{2} گھنٹے زیادہ لگتے، سفر کی شرح معلوم کیجیے۔}
\uq{234}{عرشیہ اور ابراہیم مل کر ایک کام \lr{4} گھنٹے میں مکمل کرتے ہیں۔ اگر عرشیہ اکیلی \lr{6} گھنٹے میں یہی کام مکمل کرتی ہے تو ابراہیم کو یہی کام مکمل کرنے میں کتنا وقت لگے گا؟}
\uq{235}{ایک پائپ \lr{5} گھنٹے میں جبکہ دوسرا \lr{8} گھنٹے میں پانی بھرتا ہے۔ دوسرا پائپ \lr{8} گھنٹے بعد بند کر دیا جاتا ہے کل وقت معلوم کیجیے۔}
\uq{236}{حوریہ اکیلے کام کرتے ہوئے ایک پروجیکٹ \lr{12} گھنٹے میں مکمل کر سکتی ہے۔ اگر عبد الہادی اس کی مدد کرے اور وہ دونوں مل کر اسے \lr{5} گھنٹے میں مکمل کریں، تو عبد الہادی کو اکیلے یہ پراجیکٹ مکمل کرنے میں کتنا وقت لگے گا؟}
\uq{237}{دو کاریں مخالف قصبوں یں سے ایک دوسرے کی طرف سفر کرنا شروع کرتی ہیں ان کے درمیان فاصلہ \lr{240} کلو میٹر ہے۔ ایک کار \lr{x} کلو میٹر فی گھنٹہ کی رفتار سے اور دوسری کا \lr{(10+x)} کلو میٹر کی رفتار سے سفر کرتی ہیں۔ دونوں کاریں \lr{2} گھنٹے بعد ملتی ہیں۔ ہر کار کی رفتار معلوم کیجیے۔}
\uq{238}{راشد ایک گھر کو \lr{6} دن میں رنگ کر سکتا ہے، لیکن اگر اسے مددگار مل جائے تو وہ اسے \lr{4} دن میں کر سکتا ہے۔ مددگار کو اکیلے گھر کو رنگ کرنے میں کتنا وقت لگے گا؟}
\uq{239}{اشمل \lr{10} کلو میٹر کا سفر \lr{$\dfrac{2x}{x+4}$} گھنٹوں میں کرتی ہے۔ اگر اس کی اوسط رفتار \lr{8} کلو میٹر فی گھنٹہ ہو تو \lr{x} کی قیمت معلوم کیجیے۔}
\uq{240}{فہد ایک مخصوص رفتار سے \lr{600} کلو میٹر دوڑتا ہے اور پھر اپنی رفتار دوگنی کر کے مزید \lr{600} کلو میٹر دوڑتا ہے۔ اگر اسے \lr{1200} کلو میٹر کا فاصلہ طے کرنے میں \lr{$2\dfrac{1}{2}$} گھنٹے لگیں تو وہ رفتار معلوم کیجیے جس سے اس نے دوڑنا شروع کیا۔ (میٹر فی سیکنڈ میں)}
\uq{241}{ایک سائیکل سورا ایک مخصوص رفتار سے \lr{30} کلو میٹر کا فاصلہ طے کرتا ہے۔ اگر اس کی رفتار \lr{5} کلو میٹر فی گھنٹہ تیز ہوتی، تو سفر میں ایک گھنٹہ کم لگتا۔ رفتار معلوم کیجیے۔}
\uq{242}{ایک تاجر کچھ اشیاء \lr{300} روپے میں خریدتا ہے۔ نقل و حمل کے دوران چار اشیاء ٹوٹ جاتی ہیں، لیکن وہ باقی اشیاء کو \lr{1.25} روپے فی شے منافع کے حساب سے فروخت کرتا ہے اور مجموعی طور پر \lr{25.} روپے حاصل کرتا ہے۔ اس نے کتنی اشیاء خریدی تھیں؟}
\uq{243}{ایک ٹرین خاص اوسط رفتار سے \lr{120} کلو میٹر کا فاصلہ طے کر سکتی ہے۔ سروس چیک اپ کی وجہ سے اس کی رفتار \lr{5} کلو میٹر فی گھنٹہ کم ہو جاتی ہے۔ اور اس کے نتیجے میں سفر مقرر وقت سے \lr{20} منٹ زیادہ ہو جاتا ہے۔ مقرر رفتار کیا تھی؟}
\uq{244}{مختصر کیجیے: \lr{$\dfrac{5}{x^{2}+x-6}-\dfrac{1}{2x^{2}-7x+6}$}}
\uq{245}{مختصر کیجیے: \lr{$\dfrac{x^{3}-27}{x^{2}-9}\div\dfrac{x^{2}+3x+9}{x^{2}+6x+9}$}}
\uq{246}{مختصر کیجیے: \lr{$\dfrac{a^{2}-(b-c)^{2}}{(a+b)^{2}-c^{2}}\times\dfrac{a^{2}-(b+c)^{2}}{(a-b)^{2}-c^{2}}$}}
\uq{247}{درج ذیل مساوات کو حل کیجیے: \lr{$x^{4}-16x^{2}+63=0$}}
\uq{248}{درج ذیل مساوات کو حل کیجیے: \lr{$4x^{4}-16x^{2}+15=0$}}
\uq{249}{درج ذیل مساوات کو حل کیجیے: \lr{$3^{2x}-12.3^{x}+27=0$}}
\uq{250}{درج ذیل مساوات کو حل کیجیے: \lr{$\left(x^{2}+\dfrac{1}{x^{2}}\right)-3\left(x+\dfrac{1}{x}\right)-2=0$}}
\uq{251}{درج ذیل مساوات کو حل کیجیے: \lr{$\left(x^{2}+\dfrac{1}{x^{2}}\right)+\left(x+\dfrac{1}{x}\right)=4$}}
\uq{252}{درج ذیل مساوات کو حل کیجیے: \lr{$(x+9)(x-3)(x-7)(x+5)=385$}}
\uq{253}{درج ذیل مساوات کو حل کیجیے: \lr{$(x-1)(x-2)(x+5)(x-8)=-360$}}
\uq{254}{احمد \lr{50} گلاسوں کو پینٹ کرنے میں \lr{2} گھنٹے لگاتا ہے فائزہ \lr{45} گلاسوں کو پینٹ کرنے میں \lr{2} گھنٹے لگاتی ہے۔ اکھٹے کام کرتے ہوئے انہیں \lr{150} گلاسوں کو پینٹ کرنے میں کتنا وقت لگنا چاہیے؟}
\uq{255}{ایک نل ایک ٹینک کو \lr{6} گھنٹے میں بھر سکتا ہے۔ ایک دوسرا نل اسے \lr{9} گھنٹے میں خالی کر سکتا ہے۔ اگر دونوں نل ایک ساتھ کھول دیے جائیں تو خالی ٹینک کو بھرنے میں کتنا وقت لگے گا؟}
\uq{256}{ماہم نے کھلونوں کی ایک خاص تعداد \lr{300} روپے میں خریدی، اگر ہر کھلونا \lr{5} روپے سستا ہوتا، تو وہ اس سے دو مزید کھلونے خرید سکتی تھی۔ اس نے کتنے کھلونے خریدے؟}
\uq{257}{اگر \lr{$\underline{a}=5\underline{i}-7\underline{j},\underline{b}=-\underline{i}-\underline{j}$} ہو اور \lr{$\underline{c}=2\underline{i}+3\underline{j}$} ہوں تو \lr{a+b-3c} کے متوازی یونٹ ویکٹر معلوم کیجیے۔}
\uq{258}{اگر \lr{$\underline{a}=3\underline{i}-\underline{j},\underline{b}=-2\underline{i}+4\underline{j}$} اور \lr{$\underline{c}=\underline{i}+2\underline{j}$} ہوں تو \lr{$3\underline{a}-2\underline{c}+4\underline{b}$} کے متوازی یونٹ ویکٹر معلوم کیجیے۔}
\uq{259}{دو ایسے ویکٹرز معلوم کیجیے جو \lr{$3\underline{i}-5\underline{j}$} سے دوگنا ہوں، ان میں سے ایک اس کی سمت میں ہی ہو اور دوسرا اس کی مخالف سمت میں ہو۔}
\uq{260}{اگر \lr{$\underline{a}=3\underline{b}$}، جہاں \lr{$\underline{a}=9\underline{i}-k\underline{j}$} اور \lr{$\underline{b}=3\underline{i}-12\underline{j}$} ہوں تو \lr{k} کی قیمت معلوم کیجیے۔}
\uq{261}{ثابت کیجیے کہ \lr{$\underline{a}=5\underline{i}-2\underline{j}$} اور \lr{$\underline{b}=-10\underline{i}+4\underline{j}$} متوازی ہیں۔}
\uq{262}{\lr{$\underline{a}=-5\underline{i}+2\underline{j}$} کی سمت میں یونٹ ویکٹر معلوم کیجیے۔}
\uq{263}{اگر \lr{$\underline{a}=\underline{i}+3\underline{j},\underline{c}=2\underline{i}+\underline{j}$} اور \lr{$\underline{a}+2\underline{b}=\underline{c}$} ہوں تو \lr{$|\underline{b}|$} معلوم کیجیے۔}
\uq{264}{اگر \lr{$5\underline{i}-3\underline{j}=m(\underline{i}-10\underline{j})+n(4\underline{i}-3\underline{j})$} ہوں تو \lr{m} اور \lr{n} کی قیمتیں معلوم کیجیے۔}
\uq{265}{ایک ہوائی جہاز کی ہوائی رفتار \lr{$\underline{v}_{p}=20\underline{i}$} اور کراس ونڈ \lr{$\underline{v}_{w}=50\underline{j}$} ہے۔ ہوائی جہاز کی مجموعی ولاسٹی اور اس کی عددی قیمت معلوم کیجیے۔}
\uq{266}{نقاط \lr{A(-5,4), B(-6,1)} اور \lr{C(-3,1)} کو مرتسم کر کے مثلث \lr{ABC} بنایئے۔ \lr{$\triangle ABC$} کو ٹرانسلیشن ویکٹر \lr{$8\underline{i}-6\underline{j}$} کے ذریعے \lr{$\triangle A'B'C'$} میں بھی منتقل کیجیے۔}
\uq{267}{نقاط \lr{A(-6,-2), B(-6,-5), C(-3,-5)} اور \lr{D(-3,-2)} کو مرتسم کر کے مربع \lr{ABCD} بنایئے۔ مربع \lr{ABCD} کو ٹرانسلیشن ریکٹر \lr{$9\underline{i}+7\underline{j}$} کے ذریعے مربع \lr{$A'B'C'D'$} میں بھی منتقل کیجیے۔}
\uq{268}{نقاط \lr{A(-6,3),B(-4,0), C(-2,3)} اور \lr{D(-4,4)} کو مرتسم کر کے پتنگ \lr{ABCD} بنایئے۔ پتنگ \lr{ABCD} کو ٹرانسلیشن ویکٹر \lr{$6\underline{i}-6\underline{j}$} کے ذریعے پتنگ \lr{$A'B'C'D'$} میں بھی منتقل کیجیے۔}
\uq{269}{\lr{A,B} اور \lr{D} کے محددات بالترتیب \lr{(1,2),(6,3)} اور \lr{(2,8)} ہیں، اگر \lr{ABCD} ایک توازی الاضلاع ہو تو ویکٹر کے طریقے سے \lr{C} کے محدار معلوم کیجیے۔}
\uq{270}{متوازی الاضلاع \lr{ABCD} میں دو مخالف اضلاع ویکٹرز \lr{$\overrightarrow{AB}=6\underline{i}+2\underline{j},\overrightarrow{DC}=-6\underline{i}-2\underline{j}$} کو ظاہر کرتے ہیں ویکٹر کا طریقہ سے ثابت کیجیے کے مخالف اضلاع مقدار میں برابر اور متوازی ہیں۔}
\uq{271}{نقاط \lr{A(1,2),B(4,6)} اور \lr{C(7,2)} ایک مثلث بناتے ہیں۔ ویکٹر کی عددی قیمت کو استعمال کرتے ہوئے تصدیق کیجیے کہ مثلث \lr{ABC} ایک متساوی الساقین ہے۔}
\uq{272}{ویکٹرز کا استعمال کرتے ہوئے ثابت کیجیے کہ \lr{PQRS,} ایک متوازی الاضلاع ہے، جہاں نقاط \lr{P,Q,R} اور \lr{S} کے محددات بالترتیب \lr{(1,2)(5,2)(7,6)} اور \lr{(3,6)} ہیں}
\uq{273}{ویکٹرز کا استعمال کرتے ہوئے ثابت کیجیے کے مثلث \lr{XYZ} متساوی الساقین ہے، جہاں نقاط \lr{X,Y} اور \lr{Z} کے محددات بالترتیب \lr{(0,0),(2,0)} اور \lr{(1,3)} ہیں۔}
\uq{274}{ایک گیند کو ابتدائی ولاسٹی ویکٹر \lr{$\underline{v}_{0}=10\underline{i}+20\underline{j}$} کے ساتھ پھینکا جاتا ہے۔ افقی جزو \lr{x} سمت میں ہے اور ثقلی اسراع \lr{$\underline{g}=0\underline{i}-10\underline{j}$} ہے۔ سب سے زیادہ بلندی اور افقی حد معلوم کیجیے۔}
\uq{275}{ویکٹر کے طریقے سے ثابت کیجیے کہ نقاط \lr{A(1,2), B(4,6)}، \lr{C(7,4)} اور \lr{D(4,0)} ایک متوازی الاضلاع بناتے ہیں۔}
\uq{276}{ایک پروجیکٹائل کو ابتدائی ولاسٹی ویکٹر \lr{$12\underline{i}+16\underline{j}$} کے ساتھ زمین سے فائر کیا گیا ہے۔ یہ افقی طور پر \lr{24} میٹر دور ہدف سے کتنی دیر بعد ٹکرائے گا؟}
\uq{277}{ایک ٹینک مشرق کی سمت \lr{$\underline{v}_{1}=30\underline{i}+40\underline{j}$} کی ولاسٹی سے حرکت کرتے ہوئے ایک ہدف \lr{$\underline{v}_{2}=10\underline{i}+0\underline{j}$} ولاسٹی سے گولہ فائر کرتا ہے۔ نسبتی ولاسٹی اور اس کی عددی قیمت معلوم کیجیے۔}
\uq{278}{نقاط \lr{A(-4,2), B(-1,2)} اور \lr{C(-3,4)} کو مرتسم کر کے مثلث \lr{ABC} بنایئے۔ \lr{$\triangle ABC$} کو ٹرانسلیشن ویکٹر \lr{$7\underline{i}+2\underline{j}$} کے ذریعے \lr{$\triangle A'B'C'$} میں بھی منتقل کیجیے۔}
\uq{279}{اگر \lr{A(1,2), B(4,6)} اور \lr{C(7,2)} مثلث کے رداس ہوں تو تصدیق کیجیے کہ مثلث \lr{ABC} ایک متساوی الساقین مثلث ہے۔}
\uq{280}{مثلث \lr{ABC} کو حل کیجیے، جس میں: \lr{$a=6.1\ cm,\ b=8.4\ cm,\ \alpha=42^{o}$}}
\uq{281}{مثلث \lr{abc} کو حل کیجیے، جس میں: \lr{$a=12.2\ cm,\ b=15.8\ cm,\ \gamma=50^{o}$}}
\uq{282}{مثلث \lr{ABC} کو حل کیجیے، جس میں: \lr{$b=5.2\ cm,\ c=5\ cm,\ \gamma=48^{o}$}}
\uq{283}{مثلث \lr{ABC} کو حل کیجیے، جس میں: \lr{$b=4.8\ cm,\ a=4\ cm,\ \beta=71^{o}$}}
\uq{284}{مثلث \lr{ABC} کو حل کیجیے، جس میں: \lr{$\beta=70^{o},\ b=8\ cm,\ \alpha=100^{o}$}}
\uq{285}{مثلث \lr{ABC} کو حل کیجیے، جس میں: \lr{$a=6\ cm,\ \alpha=55^{o},\ \gamma=60^{o}$}}
\uq{286}{مثلث \lr{ABC} کو حل کیجیے، جس میں: \lr{$c=7\ cm,\ \beta=34^{o},\ \gamma=64^{o}$}}
\uq{287}{مثلث \lr{ABC} کو حل کیجیے، جس میں: \lr{$b=12\ cm,\ \alpha=92^{o},\ \beta=77^{o}$}}
\uq{288}{ایک مثلث \lr{$ABC,\ \beta=45^{o},\ \gamma=30^{o},\ a=6\ cm$} ہوں تو \lr{b} اور \lr{c} معلوم کیجیے۔}
\uq{289}{مثلث \lr{$\triangle ABC,\ a=3,\ b=5,\ c=7$} ہوں تو \lr{$\alpha,\beta$} اور \lr{$\gamma$} کی قیمتیں معلوم کیجیے:}
\uq{290}{مثلث منشور میں معلوم کیجیے: \lr{$\overline{CF}$} کی لمبائی، \lr{$\overline{BF}$}، زاویہ \lr{BFC} میں جواب ایک درجہ اعشاریہ تک درست ہو۔}
\pfig{g290.png}{5.5cm}
\uq{291}{شکل ایک مثلثی مخروط کو ظاہر کرتی ہے جس کا افقی مستطیلی قاعدہ \lr{PQRS} ہے۔ جس میں \lr{$m\overline{PQ}=16\ cm,\ m\overline{QR}=10\ cm$} ہے۔ \lr{M} خط \lr{$\overline{PR}$} کا درمیانی نقطہ ہے۔ راس \lr{T} سے عمودی طور پر \lr{M} کے اوپر ہے اور \lr{$m\overline{MT}=15\ cm$} خط \lr{TP} اور قاعدہ \lr{PQRS} کے درمیان زاویہ معلوم کیجیے۔ اپنا جواب ایک درجہ اعشاریہ تک درست دیں۔}
\pfig{g291.png}{6cm}
\uq{292}{شکل مخروطی کو ظاہر کرتی ہے جس کا قاعدہ \lr{PQRS, 10} سینٹی میٹر ضلع کا ایک افقی مربع ہے۔ راس کا \lr{V} قاعدے کے وسطی نقطہ \lr{M} سے عمودی طور پر اوپر ہے اور \lr{$m\overline{VM}=12\ cm$} ہے۔ زاویہ \lr{VPM} معلوم کیجیے۔}
\pfig{g292.png}{4.5cm}
\uq{293}{\lr{ABCDE} ایک مربعی قاعدہ والا مخروط ہے۔ جس میں \lr{$m\overline{AB}=15\ cm$} اور \lr{$m\overline{AE}=m\overline{BE}=m\overline{CE}=m\overline{DC}=12\ cm$} زاویہ \lr{DEB} معلوم کیجیے اور اپنا جواب ڈگری میں دیجیے۔ (مکمل عدد)}
\pfig{g293.png}{5.5cm}
\uq{294}{دی گئی شکل ایک مثلثی مخروط کو ایک افقی مستطیلی قاعدہ \lr{ABCD} کے ساتھ ظاہر کرتی ہے۔ \lr{$m\overline{AB}=10\ cm,\ m\overline{BC}=7\ cm$} ہے۔ نقطہ \lr{M,AD} کا وسطی نقطہ ہے راس \lr{M,T} سے عمودی طور پر \lr{M} کے اوپر ہے اور \lr{$m\overline{MT}=6\ cm$} ہے۔ \lr{$\overline{TB}$} اور قاعدے کے درمیان زاویہ معلوم کیجیے۔}
\pfig{g294.png}{6.5cm}
\uq{295}{ایک آئیسومیٹرک گیم میں کیمرہ زمین سے \lr{$45^{o}$} کے زاویہ پر رکھا گیا ہے۔ اگر کھلاڑی سامنے سے \lr{10} فٹ اور اوپر \lr{10} یونٹ پر ہو تو دیکھنے کے خط کا فاصلہ معلوم کیجیے۔}
\uq{296}{ایک سرویئر ایک ٹاور کی چوٹی کو \lr{$28^{o}$} کے زاویہ صعود پر دیکھتا ہے، وہ قاعدہ سے \lr{60} میٹر کے فاصلہ پر کھڑا ہے، ٹاور کی اونچائی معلوم کیجیے۔}
\uq{297}{ایک سامع دو سپیکروں سے آنے والی آواز سنتا ہے۔ ایک سپیکر باکل سامنے \lr{6} میٹر کے فاصلے پر ہے۔ اور دوسرا \lr{$30^{o}$} کے زاویہ پر \lr{6} میٹر کے فاصلے پر ہے۔ سپیکرز کے درمیان فاصلہ معلوم کیجیے۔}
\uq{298}{ایک لائٹ ہاؤس سطح سمندر سے \lr{80} میٹر اونچی چٹان پر واقع ہے۔ ایک بحری جہاز \lr{$12^{o}$} کے زاویہ نزول پر دیکھا جاتا ہے جہاز چٹان کی بنیاد سے کتنا دور ہے؟}
\uq{299}{دی گئی شکل ایک منشور کا کراس سیکشن ہے۔ مثلث \lr{PQR} منشور کی طرف \lr{$m\overline{PR}=20\ cm,\ m\overline{MP}=12\ cm,\ m\angle PRQ=30^{o}$} اور \lr{$m\angle PQR=90^{o}$} وہ زاویہ معلوم کیجیے جو خط \lr{MR} مستوی \lr{RQLM} کے ساتھ بناتا ہے۔}
\pfig{g299.png}{5.5cm}
\uq{300}{شکل ایک مکعب نما \lr{ABCDEFGH} کو ظاہر کرتی ہے جس میں \lr{$m\overline{FC}=16\ cm,\ m\overline{AB}=8\ cm,\ m\overline{AF}=6\ cm$} ہوں تو \lr{$\overline{BC}$} کی لمبائی معلوم کیجیے۔}
\pfig{g300.png}{6cm}
\uq{301}{شکل تین مقامات \lr{A,B} اور \lr{C} کو ظاہر کرتی ہے۔ درج ذیل کا بیرنگ معلوم کیجیے: (الف) \lr{A} کا \lr{C} سے (ب) \lr{A} کا \lr{C} سے (ج) \lr{C} کا \lr{B} سے (د) \lr{B} کا \lr{C} سے}
\pfig{g301.png}{4cm}
\uq{302}{عبد الہادی \lr{100} میٹر شمال کی طرف اور پھر \lr{300} میٹر مشرق کی طرف چلتا ہے: (الف) وہ اپنی ابتدائی پوزیشن سے کتنا دور ہے؟ (ب) اپنی ابتدائی پوزیشن پر واپس جانے کے لیے اسے کس بیرنگ پر چلنا چاہیے؟}
\uq{303}{ایک بحری جہاز \lr{120} ڈگری کے بیرنگ پر \lr{20} کلو میٹر کا سفر طے کرتا ہے۔ جہاز اپنی ابتدائی پوزیشن سے کتنی دور جنوب میں گیا ہے؟}
\uq{304}{فاطمہ \lr{5.5} کلو میٹر جنوب کی طرف چلی اور پھر \lr{1.3} کلو میٹر مغرب کی طرف مڑ گئی۔ اس کے ابتدائی نقطہ سے حوریہ کا بیرنگ معلوم کیجیے۔}
\uq{305}{ایک ہوائی جہاز \lr{150} کلو میٹر مشرق کی طرف پھر \lr{100} کلو میٹر شمال مشرق (\lr{$45^{o}$} مشرق سے) کی طرف پرواز کرتا ہے۔ کل ہٹاؤ معلوم کیجیے۔}
\uq{306}{ایک بحری جہاز \lr{$045^{o}$} کے بیرنگ پر \lr{100} کلو میٹر کا سفر طے کرتا ہے پھر راستہ بدلتا ہے \lr{$135^{o}$} کے بیرنگ پر \lr{120} کلو میٹر کا سفر طے کرتا ہے۔ ابتدائی نقطہ سے آخری پوزیشن کے درمیان کا فاصلہ معلوم کیجیے۔}
\uq{307}{شکل دو بحری جہازوں \lr{P} اور \lr{Q} کی پوزیشن ظاہر کرتی ہے۔ (الف) جہاز \lr{P} اور جہاز \lr{Q} سے بیرنگ کیا ہے؟ (ب) جہاز \lr{Q} کا جہاز \lr{P} سے بیرنگ کیا ہے؟}
\pfig{g307.png}{3.8cm}
\uq{308}{تین جہاز \lr{N,M,L} شکل میں دکھائی گئی پوزیشن میں ہیں۔ جہاز \lr{M} جہاز \lr{N} کے شمال مشرق میں ہے۔ \lr{M,L} کا بیرنگ معلوم کیجیے}
\pfig{g308.png}{5cm}
\uq{309}{ایک ہوائی جہاز انج ہوائی اڈے سے روانہ ہوتا ہے اور \lr{40} کلو میٹر مشرق کی طرف سفر طے کرتا ہے جس کے بعد \lr{50} کلو میٹر شمال کی طرف چلتا ہے۔ پھر پائلٹ ممکنہ راستہ اختیار کرتے ہوئے ہوائی اڈے پر واپس آتا ہے۔ مختصر ترین راستہ کا بیرنگ اور فاصلہ کیا ہے؟}
\uq{310}{شکل میں \lr{Q} کا \lr{P} سے بیرنگ \lr{$60^{o}$} ڈگری ہے۔ \lr{P} کا \lr{Q} سے بیرنگ معلوم کیجیے۔}
\pfig{g310.png}{3.3cm}
\uq{311}{عبد اللہ ایک بحری جہاز میں نقطہ \lr{A} سے نقطہ \lr{B} تک \lr{$80^{o}$} ڈگری کے بیرنگ پر سفر کرتا ہے۔ نقطہ \lr{B} سے نقطہ \lr{A} پر واپس آنے کے لیے اسے کونسا بیرنگ لینا چاہیے؟}
\uq{312}{\lr{$\triangle ABC$} کا رقبہ معلوم کیجیے جس میں: \lr{$a=5.4\ cm,\ b=3.4\ cm,\ \alpha=49^{o}$}}
\uq{313}{\lr{$\triangle ABC$} کا رقبہ معلوم کیجیے جس میں: \lr{$\alpha=32^{o},\ \gamma=48^{o},\ c=81cm$}}
\uq{314}{ایک پائلٹ \lr{$045^{o}$} کے بیرنگ پر \lr{200} کلو میٹر پرواز کرتا ہے، پھر مڑ تا ہے اور \lr{$135^{o}$} کے بیرنگ پر \lr{150} کلو میٹر پرواز کرتا ہے۔ طیارہ اپنی ابتدائی پوزیشن سے کتنا دور ہے؟}
\uq{315}{شکل ایک مکعب نما \lr{ABCDEFGH} کو ظاہر کرتی ہے۔ \lr{$m\overline{AE}=3\ cm,\ m\overline{AB}=5\ cm,\ m\overline{BC}=7\ cm$} کی لمبائی معلوم \lr{$\overline{AG}$} کیجیے۔ اپنے جواب کو \lr{3} ہندسوں تک درست دیجیے۔}
\pfig{g315.png}{6cm}
\uq{316}{شکل ایک مکعب نما \lr{ABCDEFGH} کو ظاہر کرتی ہے۔ \lr{$m\overline{AE}=3\ cm,\ m\overline{AB}=5\ cm,\ m\overline{BC}=7\ cm$} اور \lr{$\overline{AG}$} مستوی \lr{ABCD} کے درمیان زاویہ معلوم کیجیے۔ جواب کو \lr{1} درجہ اعشاریہ تک درست دیں۔}
\pfig{g316.png}{6cm}
\uq{317}{شکل مثلثی مخروط \lr{ABCDEF} کو ظاہر کرتی ہے۔ \lr{$m\overline{AB}=8\ cm,\ m\overline{BF}=14\ cm,\ m\overline{EF}=19cm$} ہے۔ \lr{A} اور \lr{F} کے درمیان فاصلہ معلوم کیجیے۔}
\pfig{g317.png}{5.5cm}
\uq{318}{شکل مثلثی مخروط \lr{ABCDEF} کو ظاہر کرتی ہے۔ \lr{$m\overline{AB}=8\ cm,\ m\overline{BF}=14\ cm,\ \overline{AF},\ m\overline{EF}=19cm$} اور مستوی \lr{BCEF} کے درمیان زاویہ معلوم کیجیے۔}
\pfig{g318.png}{5.5cm}
\uq{319}{ہاشم \lr{4} کلو میٹر شمال کی طرف چلتا ہے پھر مڑ تا ہے اور \lr{3} کلو میٹر مشرق کی طرف چلتا ہے۔ اس کے ابتدائی نقطہ سے آخری پوزیشن کا بیرنگ معلوم کیجیے۔}
\uq{320}{ایک تعمیر میں تین سٹیل کی سلاخیں نقاط \lr{A,B} اور \lr{C} (خط مستقیم میں نہیں) پر نصب ہیں۔ ان تینوں نقاط میں سے کتنے چھلے گزر سکتے ہیں۔}
\uq{321}{ایک پارک میں روشنی کے کھمبے \lr{10m} دائروی ٹکڑے کے کنارے پر رداس کے \lr{14m} فاصلہ پر نصب ہیں جیسا کہ شکل میں دکھایا گیا ہے۔ پارک کے مرکز سے وتر کا فاصلہ معلوم کیجیے۔}
\pfig{g321.png}{4.5cm}
\uq{322}{ایک دائرے میں دونوں وتروں \lr{AB} اور \lr{CD} کی لمبائی \lr{10cm} ہے۔ اگر مرکز سے \lr{$\overline{AB}$} کا فاصلہ \lr{6cm} ہے تو مرکز سے \lr{$\overline{CD}$} کا فاصلہ کیا ہے؟}
\uq{323}{ایک \lr{10 cm} رداس کے دائروی میز کی بالائی سطح پر \lr{12cm} کے فاصلے پر دو سوراخ \lr{A} اور \lr{B} کیے گئے ہیں۔ مرکز سے \lr{$\overline{AB}$} کا عمودی فاصلہ معلوم کیجیے۔}
\uq{324}{دائرہ جس کا مرکز \lr{O} ہے۔ وتر \lr{PQ} اور \lr{RS} کا مرکز سے عمودی فاصلہ \lr{6cm} ہے۔ اگر وتر \lr{PQ} کی لمبائی \lr{18cm} ہو تو دوسرے وتر کی لمبائی معلوم کیجیے۔}
\uq{325}{دائرے کے مرکز سے ایک خط قائمہ زاویہ پر \lr{10cm} کے وتر کو قطع کرتا ہے یہاں دائرہ کا رداس \lr{6cm} ہے۔ مرکز سے وتر کا فاصلہ کیا ہے؟}
\uq{326}{ایک دائرے جس کے مرکز سے وتر \lr{AB} پر عمود کھینچا گیا ہے۔ اگر \lr{$m\overline{AB}=12\ cm$} ہو تو تنصیب کرنے کے بعد ہر حصے کی لمبائی کیا ہے؟}
\uq{327}{ایک دائرہ میں \lr{$m\overline{AB}=8\ cm$} کو مرکز \lr{O} سے ایک خط نقطہ \lr{P} پر دو برابر حصوں میں تقسیم کرتا ہے۔ اگر \lr{A} سے \lr{P} کا فاصلہ \lr{4cm} ہے اور رداس \lr{5cm} ہے تو \lr{$m\overline{OP}$} معلوم کریں۔}
\uq{328}{\lr{13} میٹر رداس کے دائروی محراب میں ایک \lr{24} میٹر کا وتر \lr{AB} ہے۔ مرکز سے وتر تک اونچائی معلوم کریں۔}
\uq{329}{ایک وتر کی لمبائی معلوم کیجیے جو ایک دائرے کے مرکز سے \lr{5cm} کے فاصلے پر ہے، جس کا رداس \lr{13cm} ہے۔}
\uq{330}{ایک دائروی ٹریک میں دو متماثل قوسیں مرکز پر \lr{$60^{o}$} کے زاویے بناتی ہیں۔ اگر دائرے کے ایک وتر کی لمبائی \lr{10m} ہو تو دوسرے وتر کی لمبائی معلوم کیجیے۔}
\uq{331}{ایک دائروی فوارے میں پانی کے دو تیز دھارے اس طرح نصب کیے گئے ہیں کہ وہ برابر لمبائی والی قوسوں کے ساتھ پانی کا چھڑ کاؤ کرتے ہیں۔ اگر ایک دھارا \lr{P} اور \lr{Q} کے درمیان چھڑ کاؤ کرے اور خط مستقیم کا فاصلہ (وتر \lr{PQ}) \lr{12m} ہو تو دوسرے تیز دھارے کی طرف سے طے کیے گئے خط مستقیم (وتر \lr{RS}) کا فاصلہ کیا ہے جو قوم \lr{RS} کے ساتھ سپرے کرتا ہے جو کہ قوس \lr{PQ} کے متماثل ہے؟}
\uq{332}{ایک دائروی پارک میں دور راستے \lr{$\overline{AB}$} اور \lr{$\overline{CD}$} سیدھے راستے (وتر) ہیں جن کی لمبائی \lr{14m} ہے۔ ان راستوں سے بننے والی قوس صغیرہ کے بارے میں کیا کہا جاسکتا ہے؟}
\uq{333}{دو متماثل دائروی گھڑیوں میں منٹ والی سوئی دونوں گھڑیوں میں مرکز سے \lr{3} کی طرف اشارہ کرتی ہے۔ ایک آرائشی تار دونوں گھڑیوں پر \lr{3} کو \lr{9} سے ملاتی ہے۔ کیا دونوں گھڑیوں پر \lr{3} سے \lr{9} تک قوسیں متماثل ہیں؟}
\uq{334}{دائروی باغ میں دو سجاوٹی لائنیں اس طرح رکھی گئی ہیں کہ قوسوں کے ہر جوڑے کے درمیان \lr{(A,B,C)} اور \lr{D}) کی لمبائی برابر ہے۔ اگر نقاط \lr{A} اور \lr{B} کے درمیان خط کا سیدھا فاصلہ \lr{8m} ہو تو نقاط \lr{C} اور \lr{D} کے درمیان فاصلہ کیا ہے؟}
\uq{335}{ایک دائروی پارک میں دو بنچ رکھے گئے ہیں اور ان کی پوزیشنوں سے بننے والے وتروں کی لمبائی \lr{10m} ہے۔ آپ ہر بنچ سے مرکز پر بننے والے زاویے کے بارے میں کیا کہہ سکتے ہیں؟}
\uq{336}{جنت ایک مثلثی باغ کا نمونہ بنا رہی ہے جس کے رداس \lr{A(2,1), B(4,4)} اور \lr{C(1,5)} ہیں۔ کیا وہ ان تینوں راسوں کو چھونے والا ایک دائروی فوراہ نصب کر سکتی ہے؟}
\uq{337}{دو وتر \lr{PQ} اور \lr{RS}، ہر ایک لمبائی \lr{18cm} ہے، اگر وتر \lr{PQ} کا دائرے کے مرکز سے فاصلہ \lr{7cm} ہو تو وتر \lr{RS} کا مرکز سے فاصلہ معلوم کیجیے۔}
\uq{338}{ایک درخت کی شاخ ایک دائروی تالاب کے پار پھیلی ہوئی ہے، جو \lr{20m} کا وتر ہوتی ہے۔ ایک پیمائشی سلاخ جس کی لمبائی \lr{10m} ہے وتر پر مرکز سے عمود ہے، تالاب کا رداس کیا ہے؟}
\uq{339}{ایک سٹیل کی سلاخ \lr{6m} لمبی ہے، جو ایک گول ڈھانچے کے اندرونی حصے میں وتر کے طور پر پڑی ہے۔ مرکز سے ایک عمودی سلاخ اس کی تنصیف کرتی ہے۔ اگر مرکز سے سلاخ کا عمودی فاصلہ \lr{4m} ہو تو دائروی ڈھانچے کا رداس معلوم کیجیے۔}
\uq{340}{ایک دائروی کلاک جس کا رداس \lr{6cm} ہے، \lr{2} اور \lr{10} کے نقاط ایک وتر بناتے ہیں جس کی لمبائی \lr{10cm} ہے۔ کلاک کے مرکز سے وتر کا عمودی فاصلہ معلوم کیجیے۔}
\uq{341}{ایک مشین میں مناسب گردش کے لیے دو دائروی گیئرز ایک دوسرے کو بیرونی طور پر مس کرتے ہیں۔ دو دائروی گیئرز کے رداس \lr{5cm} اور \lr{7cm} ہیں۔ اگر وہ بیرونی طور پر مس کرتے ہیں تو ان کے مراکز کا درمیانی فاصلہ کیا ہے؟}
\uq{342}{ایک چھوٹا سینسر ایک سیٹلائٹ ڈش کے اندر واقع ہے اور اس کی دیوار کو اندرونی طور پر مس کرتا ہے۔ اگر ڈش کا رداس \lr{15cm} ہو اور سینسر کا رداس \lr{2.5cm} ہو تو ان کے مراکز کا درمیانی فاصلہ معلوم کریں۔}
\uq{343}{ایک اندرونی ہولڈر بیرونی سلنڈر نما برتن کو اندرونی طور پر مس کرتا ہے۔ اگر بیرونی برتن کا رداس \lr{8cm} ہو اور اندرونی پولڈر کا رداس \lr{6cm} ہو تو ان کے مراکز کا درمیانی فاصلہ معلوم کیجیے۔}
\uq{344}{\lr{A,B} اور \lr{C} مرکز \lr{O} کے ساتھ ایک دائرہ کے محیط پر نقاط ہیں۔ \lr{$\overleftrightarrow{DE}$} اور \lr{$\overline{AC}$} دائرہ کا قطر ہے اور نقطہ \lr{C} پر دائرہ کا مماس اور \lr{$m\angle BCD=62^{o}$} ہے۔ معلوم کیجیے۔ (الف) \lr{$m\angle BCA$} (ب) \lr{$m\angle BAC$}}
\uq{345}{ایک تاریخی یاد گار میں \lr{3m} رداس کا ایک دائروی فوراہ بنایا گیا ہے۔ فوارے کے مرکز سے \lr{7m} کے فاصلے پر ایک پرچم کا پول نصب ہے۔ پول سے دو رسیاں فوارے کے کنارے سے اس طرح باندھی گئی ہیں کہ وہ صرف اس مس کرتی ہیں۔ ہر رسی کی لمبائی معلوم کیجیے۔}
\uq{346}{ایک اہرام کی شکل کا مجسمہ \lr{10m} رداس کے دائروی پلازہ کے مرکز میں رکھا گیا ہے۔ ایک آرائشی پول دائرہ کے مرکز سے \lr{26m} کے فاصلے پر نصب ہے۔ پول کے بالائی حصے سے پلازہ کے کنارے تک دو گائیڈ تاریں جوڑی گئی ہیں، جو صرف دائرہ کو مس کرتی ہیں۔ ہر تار کی لمبائی معلوم کیجیے۔}
\uq{347}{مرکز \lr{P} اور رداس \lr{4cm} کے دائرہ کا نقطہ \lr{Q} پر مماس \lr{$\overline{QR}$} ہے۔ یہ قطعہ خط \lr{PR} سے اس طرح ملتا ہے کہ \lr{$\overline{QR}-m\overline{PR}=8\ cm$} کی لمبائی کیا ہے؟}
\uq{348}{ایک ماہر تعمیرات ایک عمارت پر ایک دائروی گنبد بناتا ہے جس کا رداس \lr{5m} ہے۔ گنبد کے مرکز سے روشنی کا راڈ \lr{12m} دور نصب کیا گیا ہے۔ سہارا دینے والی تاریں راڈ سے گنبد کے ساتھ منسلک کی گئی ہیں، جو گنبد کو دو نقاط پر چھو ہی ہیں۔ ہر تار کی لمبائی معلوم کریں۔}
\uq{349}{ایک خم دار پٹڑی ایک قوس بناتی ہے جو دائروی میدان کے مرکز پر \lr{$60^{o}$} کا زاویہ بناتی ہے۔ یہ حد کسی ایک نقطہ پر کیا زاویہ بناتی ہے؟}
\uq{350}{\lr{A,B,C} اور \lr{D} مرکز \lr{O} والے ایک دائرے کے محیط پر نقاط ہیں جیسا کہ شکل میں دکھایا گیا ہے۔ زاویہ \lr{x} معلوم کیجیے۔ (تصویر)}
\uq{351}{ایک دائروی محراب پر ایک یاد گار کے داخلی دروازے کے اوپر، دو مرتکز روشنی کے لیمپ اس طرح لگائے گئے ہیں کہ وہ ایک ہی وتر \lr{PQ} سے محراب پر دو مخالف نقاط پر روشنی ڈالتے ہیں۔ اگر ایک لیمپ پر بننے والا زاویہ \lr{$55^{o}$} ہو تو \lr{PQ} کے اسی طرف کے دوسرے لیمپ پر بننے والا زاویہ معلوم کیجیے۔}
\uq{352}{ایک دائروی جھولے کا رداس \lr{12m} ہے۔ اگر ایک مسافر جھولے کے محیط کے ساتھ \lr{18m} کا فاصلہ طے کرتا ہے تو مسافر کے مقام کا نقطہ آغاز سے زاویہ (ریڈین) معلوم کیجیے۔}
\uq{353}{ملحقہ شکل میں \lr{$m\angle BDC=74^{o}$} ہے۔ \lr{$m\angle BAC,\ m\angle BOC$} اور \lr{$m\angle OBC$} معلوم کیجیے۔ (تصویر)}
\uq{354}{ایک سیکٹر کا محیط اور رقبہ بالترتیب \lr{14cm} اور \lr{$10\ cm^{2}$} ہیں۔ دائرہ کا رداس اور سیکٹر کا مرکزی زاویہ معلوم کیجیے۔}
\uq{355}{ایک لائبریری کی بڑی دائروی رنگین شیشے کی کھڑکی میں ایک ہی وتر \lr{AB} سے دو آرائشی بیم کے مخالف قوس \lr{AB} پر دو نقاط \lr{P} اور \lr{Q} تک کھینچے گئے ہیں اس طرح کہ \lr{$\angle APB$} اور \lr{$\angle AQB$} بنتے ہیں۔ اگر \lr{$\angle APB$} کی پیمائش \lr{$40^{o}$} ہو تو \lr{$\angle AQB$} کی پیمائش کیا ہے؟}
\uq{356}{شکل میں \lr{P,Q} اور \lr{R} ایک دائرے کے محیط پر نقاط ہیں اور قطعہ خط \lr{PQ} قطر ہے۔ \lr{x} اور \lr{y} زاویے معلوم کیجیے۔}
\uq{357}{دو گول کھانے کی میزیں اپنے کناروں کو مس کر رہی ہیں۔ ان میزوں کے رداس \lr{10cm} اور \lr{12cm} ہیں۔ اگر وہ بیرونی طور پر مس کرتی ہوں تو ان کے مراکز کا درمیانی فاصلہ معلوم کیجیے۔}
\uq{358}{ایک دائروی پچی کاری نمونہ ایک جدید اہرام نما عمارت کے فرش پر رکھا گیا ہے۔ دو ٹائلیں اس طرح رکھی گئی ہیں کہ وہ ایک ہی وتر \lr{EF} سے مخالف قوس پر دو نقاط \lr{G} اور \lr{H} تک پہنچتی ہیں۔ اگر \lr{$m\angle EGF=65^{o}$} ہو تو \lr{$m\angle EHF$} معلوم کیجیے۔}
\uq{359}{ایک دائروی سالگرہ کیک کا رداس \lr{15cm} ہے۔ ایک ٹکڑے کی قوس کی لمبائی \lr{10cm} لمبی ہے۔ اس ٹکڑے کا رقبہ معلوم کیجیے۔}
\uq{360}{ایک دائرے کے \lr{5cm} رداس سے قطع کیے ہوئے ایک سیکٹر کا محیط \lr{16cm} ہے۔ اس سیکٹر کا رقبہ معلوم کیجیے۔}
\uq{361}{کوئی سے تیم غیر ہم خط نقاط (پارک میں موجود روشنی کے کھمبوں کی جگہوں پر) لیں اور ان تمام نقاط سے گزرنے والا دائرہ بنائیں۔}
\uq{362}{\lr{ABC} دائروی فوارے کی قوس ہے، مرکز معلوم کیے بغیر دائرہ مکمل کیجیے۔}
\uq{363}{دائروی جھولے کے رم کا ایک حصہ ایک قوس کے طور پر دکھائی دے رہا ہے۔ قوس پر کوئی سے تین نقاط کا استعمال کرتے ہوئے، مرکز معلوم کر کے دائرہ بنائیں۔}
\uq{364}{نقطہ \lr{P} سے قوس \lr{$\overset{\frown}{ABC}$} کا مماس کھینچیں، جب \lr{P} قوس کے باہر ہو۔}
\uq{365}{\lr{1.6cm} رداس کا ایک دائرہ کھینچیں اور دو مماس کھینچیں جو \lr{$30^{o}$} کے زاویہ پر ملیں۔}
\uq{366}{ایک پلی سے \lr{r=2.1cm} رداس والی ایک مشین کے پہیے پر دو مماس بنائیں، اس طرح کہ ان کا درمیانی زاویہ \lr{$30^{o}$} ہو۔}
\uq{367}{\lr{1.7cm} رداس کا ایک دائرہ بنایئے اور دو مماس کھینچیں جو \lr{$30^{o}$} کے زاویہ پر ملتے ہوں۔}
\uq{368}{دائروی ٹریک کا ایک حصہ لیں۔ اس کی قوس کا مرکز معلوم کرنے کے لیے وتروں اور عمودی ناصفوں کا استعمال کیجیے۔}
\uq{369}{دو آرائشی باریں ایک دائروی پھولوں کے تختے کو مس کرتی ہیں جس کا رداس \lr{2.1cm} ہے اور یہ باہر کی طرف \lr{$30^{o}$} کے زاویہ پر مل کر ایک داخلی محراب بناتی ہے۔ اس نقطہ سے جہاں یہ \lr{$30^{o}$} کے زاویہ پر ملتی ہیں۔ پھولوں کے تختے پر دو مماس کھینچیں۔}
\uq{370}{درج ذیل تعددی تقسیم کے ایک گروپ کی یاداشت کے ٹیسٹ پر کارکردگی کو ظاہر کرتی ہے: (الف) گروپ سکور کا مجموعی تعدد کثیر الاضلاع گراف بنائیں۔ (ب) وسطانیہ، زیریں اور بالائی چہاری حصے \lr{$D_{3},P_{90}$} اور اندرونی چہاری حصے کی سعت گراف کی مدد سے معلوم کیجیے۔ (نوٹ: تمام جوابات مکمل اعداد میں لکھیں)}
\dtab{|c|c|c|c|c|c|c|}{69-73&64-68&59-63&54-58&49-53&44-48&\urdu{\textbf{سکور}}\\ \hline 8&12&23&15&12&3&\urdu{\textbf{طلبہ کی تعداد}}}
\uq{371}{درج ذیل سکور کی تقسیم کے لیے ترسیم بنایئے:}
\dtab{|c|c|c|c|c|c|}{109-119&99-109&89-99&79-89&69-79&\urdu{\textbf{سکور}}\\ \hline 5&8&13&10&4&\urdu{\textbf{f}}}
\usc{وسطانیہ زیریں اور بالائی چہاری حصے \lr{$D_{6},P_{80}$} اور اندرونی چہاری حصہ کی سعت گراف کی مدد سے معلوم کیجیے۔ (نوٹ: تمام جوابات مکمل اعداد میں لکھیں)}
\uq{372}{دیے گئے مواد کے لیے \lr{$Q_{1},Q_{3}$} وسطانیہ، سعت اور اندرونی چہاری حصے کی سعت معلوم کیجیے۔ \lr{$3,6,8,4,7,5,10,11,13,9,14,12,15,16,17$} باکس اور وسکر پلاٹ بھی بنایئے:}
\uq{373}{درج ذیل مواد کے لیے \lr{$Q_{1},Q_{3}$} وسطانیہ، سعت، اندرونی چہاری حصے کی سعت اور انتہائی قدروں کو معلوم کیجیے: \lr{$102,98,95,100,93,110,108,104,97,96,101,99,105,107$} باکس اور وسکر پلاٹ بھی بنایئے۔}
\uq{374}{درج ذیل جدول میں فٹ بال کھلاڑیوں کے \lr{200} میں سے ٹیسٹ سکورز دیے گئے ہیں (ایک رفتار کی پیمائش اور دوسرا مضبوطی کی پیمائش)}
\dtab{|c|c|c|c|c|c|c|c|c|c|}{150&136&50&110&100&40&70&38&20&\urdu{\textbf{پہلا ٹیسٹ}}\\ \hline 36&44&140&50&70&120&130&142&160&\urdu{\textbf{دوسرا ٹیسٹ}}}
\usc{(i) انتشاری خاکہ بنائے اور موزوں ترین خط کھنچیں (ii) پہلے اور دوسرے ٹیسٹ کے درمیان ارتباط کی وضاحت کیجیے (iii) عبد اللہ نے دوسرے ٹیسٹ میں \lr{60} سکور کیا۔ اس کا پہلے ٹیسٹ میں سکور معلوم کیجیے۔}
\uq{375}{درج ذیل جدول میں یونیورسٹی کے طلبہ کے روزانہ سوشل میڈیا کے استعمال کے گھنٹے اور \lr{GPA} دیے گئے ہیں:}
\dtab{|c|c|c|c|c|c|c|c|c|c|c|}{9&8&7&6&5&4&3&2&1&0.5&\urdu{\textbf{سوشل میڈیا کے گھنٹے}}\\ \hline 2.5&2.6&2.8&2.9&3.1&3.3&3.5&3.7&3.9&4.0&\urdu{\textbf{GPA}}}
\usc{(i) انتشاری خالہ بنایئے اور موزوں ترین خط کھینچیں (ii) سوشل میڈیا کے استعمال اور \lr{GPA} کے درمیان ارتباط کی شناخت کیجیے (iii) اس طالب علم کا \lr{GPA} معلوم کیجیے جو سوشل میڈیا پر \lr{6.5} گھنٹے صرف کرتا ہو۔}
\uq{376}{کرن نے پچھلے سال کے دوران ہر ماہ بارش یا برف والے دنوں کی تعداد اور ریکارڈ کی درج ذیل باکس اور وسکر پلاٹ اس کے مواد کو ظاہر کرتا ہے۔ (تصویر) (الف) انتہائی قدریں بارش یا برف باری والے دنوں کی تعداد کے بارے میں کیا بتاتی ہیں؟ (ب) بارش یا برف باری والے دنوں کا وسطانیہ معلوم کیجیے۔ (ج) زیریں اور بالائی چہاری حصے کیا ہیں؟ (د) کہاں زیادہ تر مواد مرتکز ہے؟}
\uq{377}{کالج کے \lr{50} طلبہ کے \lr{100} میں سے حاصل کردہ مارکس دیے گئے ہیں۔ مجموعی تعددی جدول بنایئے۔}
\dtab{|c|c|c|c|c|c|c|c|c|c|}{74&71&64&77&91&66&45&76&66&85\\ \hline 78&80&67&71&58&70&42&76&78&47\\ \hline 84&78&69&65&72&81&87&68&48&73\\ \hline 83&73&92&62&72&56&87&58&53&75\\ \hline 74&88&79&62&72&58&61&51&81&97}
\uq{378}{جدول ایک سپیڈ چیک پوائنٹ سے گزرنے والی \lr{90} گاڑیوں کی رفتار (کلو میٹر فی گھنٹہ) کے بارے میں معلومات فراہم کرتا ہے۔ کمی کے طریقہ سے مجموعی تعددی کثیر الاضلاع کی تشکیل کریں اور وسطانیہ، \lr{$Q_{1},Q_{3},D_{4},P_{95}$} اور اندرونی چہاری حصہ کی سعت (\lr{IQR}) کی قیمتیں معلوم کیجیے۔}
\dtab{|c|c|c|c|c|c|}{80-90&70-80&60-70&50-60&40-50&\urdu{\textbf{رفتار (کلو میٹر فی گھنٹہ)}}\\ \hline 9&17&30&18&10&\urdu{\textbf{f}}}
\uq{379}{درج ذیل تعددی تقسیم کے لیے مجموعی تعددی کثیر الاضلاع کو زیادہ کے طریقے سے کھینچیے۔ جو مزدوروں کی تعداد کی ہفتہ وار اجرتیں (روپے ہزاروں میں) ظاہر کرتا ہے اور اسی گراف کی مدد سے \lr{$Q_{2}$} بھی معلوم کریں۔}
\dtab{|c|c|c|c|c|c|}{80-100&60-80&40-60&20-40&0-20&\urdu{\textbf{ہفتہ وار اجرتیں (روپے ہزاروں میں)}}\\ \hline 21&40&61&48&30&\urdu{\textbf{مزدوروں کی تعداد}}}
\uq{380}{نیچے دیا گیا جدول دس طلبہ کے دہرائی میں صرف شدہ وقت اور اس کے متعلقہ ٹیسٹ سکور کو ظاہر کرتا ہے۔}
\dtab{|c|c|c|c|c|c|c|c|c|}{5&7&3&2&6&4&1&0.5&\urdu{\textbf{دہرائی پر صرف کیا گیا وقت (گھنٹوں میں)}}\\ \hline 60&70&46&30&68&62&38&22&\urdu{\textbf{ٹیسٹ سکور (\%)}}}
\usc{(الف) انتشاری خاکہ بنائیں اور موزوں ترین خط کھینچیں (ب) انتشاری خاکہ میں دکھائی گئی ارتباط کی وضاحت کیجیے (ج) سلیم نے دہرائی ٹیسٹ پر \lr{5.5} گھنٹے صرف کیے۔ موزوں ترین خط کا استعمال کرتے ہوئے ٹیسٹ کا سکور معلوم کیجیے۔}
\uq{381}{\lr{7} کارکنوں کے ایک گروہ نے ایک دن میں تیار کی جانے والی اشیا کی تعداد بتائی: \lr{52،55،50،53،54،56،52} تیار کی گئی اشیاء کا معیاری انحراف اور تغیریت معلوم کیجیے۔}
\uq{382}{ایک لائبریرین نے ہفتے کے \lr{5} دنوں میں آنے والے افراد کی تعداد ریکارڈ کی: \lr{120،135،130،125،140} آنے والے افراد کی تغیریت اور معیاری انحراف معلوم کیجیے۔}
\uq{383}{پہلے \lr{23} طاق اعداد کی سعت تغیریت اور معیاری انحراف معلوم کیجیے۔}
\uq{384}{ایک ہفتے میں پانچ اضلاع کے مختلف مقامات پر بارش کا ریکارڈ ذیل ہے اس کی تغیریت اور معیاری انحراف معلوم کیجیے:}
\dtab{|c|c|c|c|c|c|}{63&61&54&51&42&\urdu{\textbf{بارش (ملی میٹر)}}\\ \hline 5&9&4&13&5&\urdu{\textbf{مقامات کی تعداد}}}
\uq{385}{مشین \lr{A} کی آؤٹ پٹ (یونٹس): \lr{98،100،102،101،99}، مشین \lr{B} کی آؤٹ پٹ (یونٹس): \lr{95،100،105،90،110} (الف) کس مشین نے بہتر کارکردگی دکھائی؟ (ب) کون سی مشین کا تسلسل برقرار رہا؟}
\uq{386}{دو سیلز پرسنز (اشخاص) کی ماہانہ فروخت (روپے لاکھوں میں) \lr{6} مہینوں کے دوران درج ذیل ہے: شخص \lr{A: 5.5،5.7،5.4،5.6،5.8،5.6} شخص \lr{B: 5.5،6.5،4.5،6.0،5.0،6.0} ان کی کارکردگی اور مستقل مزاجی کا موازنہ کیجیے۔}
\uq{387}{دیا گیا جدول ایک ٹیکسٹائل مل کے کارکنوں کی یومیہ اجرتوں کو ظاہر کرتا ہے جسے چھ آمدنی کے گروہوں میں تقسیم کیا گیا ہے۔}
\dtab{|c|c|c|c|c|c|c|}{1800-2000&1600-1800&1400-1600&1200-1400&1000-1200&800-1000&\urdu{\textbf{یومیہ اجرت (روپے)}}\\ \hline 1&2&8&6&4&2&\urdu{\textbf{تعدد}}}
\usc{اجرتوں کی اوسط تغیریت اور معیاری انحراف معلوم کیجیے۔}
\uq{388}{پانچ صوبوں میں بے روزگاری کی شرح (فیصد میں) \lr{5.2،6.0،4.8،5.5} اور \lr{6.2} ہیں۔ معیاری انحراف معلوم کیجیے اور یہ بیان کیجیے کہ کیا بے روزگاری کی شرح متوازن ہے؟}
\uq{389}{تغیریت اور معیاری انحراف معلوم کیجیے: \lr{$\Sigma x=45,\ \Sigma x^{2}=421,\ n=5$}}
\uq{390}{تغیریت اور معیاری انحراف معلوم کیجیے: \lr{$\Sigma fx=210,\ \Sigma fx^{2}=7560,\ \Sigma f=6$}}
\uq{391}{تغیریت اور معیاری انحراف معلوم کیجیے: \lr{$\bar{x}=18,\ \Sigma fx^{2}=1670,\ \Sigma f=5$}}
\uq{392}{دیا گیا جدول ملازمین کے ایک گروپ کی ہفتہ وار آمدنی (روپے ہزاروں میں) کو ظاہر کرتا ہے۔ ہفتہ وار آمدنی کی سعت معلوم کیجیے۔}
\dtab{|c|c|}{\urdu{\textbf{ہفتہ وار آمدنی (ہزار روپے میں)}}&\urdu{\textbf{ملازمین کی تعداد}}\\ \hline 25-35&4\\ 35-45&6\\ 45-55&5\\ 55-65&3}
\uq{393}{ایک فیکٹری درج ذیل لمبائیوں (سینٹی میٹر) کی سلاخیں تیار کرتی ہے: \lr{50.1،49.9،50.0،50.2،49.8} معیاری انحراف کا استعمال کرتے ہوئے یکسانیت کی جانچ کیجیے۔}
\uq{394}{ایک مالیاتی تجزیہ کار نے ایک مخصوص دن چند کمپنیوں کے حصص کی قیمتوں کے اختتامی سکورز کو درج ذیل طریقے گروہوں میں ترتیب دیا:}
\dtab{|c|c|c|c|c|c|}{80-90&70-80&60-70&50-60&40-50&\urdu{\textbf{حصص کی قیمتیں (روپے)}}\\ \hline 2&4&6&3&2&\urdu{\textbf{تعدادات}}}
\usc{اوسط، تغیریت اور معیاری انحراف معلوم کیجیے۔ آپ حصص کی قیمتوں کے اتار چڑھاؤ کے بارے میں کیا نتیجہ اخذ کر سکتے ہیں؟}
\uq{395}{دو طلبہ، علی اور روما نے ریاضی کے \lr{5} ٹیسٹوں میں حصہ لیا۔ ان کے سکورز (\lr{100} میں سے) درج ذیل تھے: علی: \lr{72،68،74،70،76} روما: \lr{75،68،80،65،72} (الف) اوسطاً کس نے بہتر کارکردگی دکھائی؟ (ب) کون زیادہ مستقل مزاج تھا۔}
\uq{396}{\lr{5} اعداد کا مجموعہ \lr{45} ہے اور ان کے مربعوں کا مجموعہ \lr{421} ہے۔ مواد کے سیٹ کا اوسط اور معیاری انحراف معلوم کیجیے۔}
\uq{397}{لمبی چھلانگ کے لیے درج ذیل نتائج ریکارڈ کیے گئے:}
\dtab{|c|c|c|c|c|c|c|}{220-230&210-220&200-210&190-200&180-190&170-180&\urdu{\textbf{فاصلہ (سینٹی میٹروں میں)}}\\ \hline 3&8&12&9&6&2&\urdu{\textbf{f}}}
\usc{مجموعی تعدد کثیر الاضلاع بنایئں اور وسطانیہ \lr{$Q_{1},Q_{3},D_{7},P_{90}$} اور اندرونی چہاری حصے کی سعت بھی معلوم کیجیے۔}
\uq{398}{درج ذیل جدول دس افراد کے لیے ہفتہ وار ورزش کے گھنٹے اور آرام کے وقت دل کی دھڑکن فی منٹ کو ظاہر کرتا ہے:}
\dtab{|c|c|c|c|c|c|c|c|c|c|}{8&7&6&5&4&3&2&1&0&\urdu{\textbf{ورزش کے گھنٹے}}\\ \hline 68&70&72&74&76&80&83&85&90&\urdu{\textbf{دل کی دھڑکن}}}
\usc{(الف) مواد کا انتشاری خاکہ بنایئے اور موزوں ترین خط کھینچیں (ب) مشاہدہ کیے گئے ارتباط کی قسم بیان کیجیے (ج) سکینہ ہر ہفتے \lr{6.5} گھنٹے ورزش کرتی ہے اس کے دل کی دھڑکن معلوم کیجیے۔}
\uq{399}{درج ذیل مواد کی سعت، تغیریت اور معیاری انحراف معلوم کیجیے:}
\dtab{|c|c|c|c|c|c|c|}{25-30&20-25&15-20&10-15&5-10&0-5&\urdu{\textbf{جماعتی وقفہ}}\\ \hline 4&16&24&20&12&4&\urdu{\textbf{f}}}
\uq{400}{\lr{6} مہینوں کے لیے دو خاندانوں کے ماہانہ گھریلو اخراجات (روپے ہزاروں میں) درج ذیل تھے۔ خاندان \lr{A: 45،47،46،48،46،47} ، خاندان \lr{B: 38،52،40،50،42،49} ماہانہ اخراجات کا اوسط اور معیاری انحراف معلوم کیجیے۔ کون سا خاندان اوسطاً زیادہ خرچ کرتا ہے؟ کس خاندان کے اخراجات زیادہ مستحکم ہیں؟}
\uq{401}{ایک فیکٹری میں \lr{40} ملازموں کی یومیہ اجرتیں درج ذیل جدول میں دی گئی ہیں:}
\dtab{|c|c|c|c|c|c|}{1800-2000&1600-1800&1400-1600&1200-1400&1000-1200&\urdu{\textbf{یومیہ اجرت (روپے)}}\\ \hline 4&8&12&10&6&\urdu{\textbf{ملازموں کی تعداد}}}
\usc{یومیہ اجرت کا وسط، تغیریت اور معیاری انحراف معلوم کیجیے۔}
\uq{402}{ایک استاد بلا ترتیب \lr{2} لڑکوں (حامد، احمد) اور دو لڑکیوں (زینب، سمعیہ) کے ایک گروہ سے ایک لڑکا اور ایک لڑکی کا انتخاب کرتا ہے۔ شجری خاکہ بنایئے اور تمام ممکنہ نتائج کے لیے سیمپل اسپیس کی فہرست بنایئے۔}
\uq{403}{ایک ڈائس اور ایک سکہ کو ایک ساتھ اچھالا گیا: (الف) امکانی خاکہ بنا کر سیمپل اسپیس معلوم کیجیے۔ (ب) شجری خاکہ کے ذریعے سیمپل اسپیس معلوم کیجیے۔ (ج) ایک ٹیل اور ایک جفت عدد آنے کا احتمال کیا ہے؟}
\uq{404}{دو غیر جانبدار ڈائس کو ایک بار اچھالا گیا: (الف) شجری خاکہ کے ذریعے سیمپل اسپیس معلوم کیجیے۔ (ب) امکانی خاکہ کے ذریعے سیمپل اسپیس معلوم کیجیے۔ \lr{(a)} دونوں ڈائس پر ایک ہی عدد ہو۔ \lr{(b)} ان کا حاصل ضرب ایک مفرد عدد ہے۔ \lr{(c)} مجموعہ ایک جفت عدد ہو۔ \lr{(d)} مجموعہ \lr{13} ہو۔}
\uq{405}{تین غیر جانبدار سکوں کو ایک ساتھ اچھالا گیا۔ درج ذیل حاصل کرنے کا احتمال معلوم کیجیے: (الف) تمام ٹیل (ب) کم از کم ایک ہیڈ (ج) زیادہ سے زیادہ دو ٹیل}
\uq{406}{تین غیر جانبدار سکوں کو ایک ساتھ اچھالا گیا۔ درج ذیل حاصل کرنے کا احتمال معلوم کیجیے: (الف) دو ہیڈز (ب) زیادہ سے زیادہ دو ہیڈز (ج) کوئی ہیڈ نہیں}
\uq{407}{ایک تھیلے میں \lr{4} سرخ گیندیں، \lr{5} سفید گیندیں، \lr{6} سبز گیندیں اور \lr{3} سیاہ گیندیں ہیں۔ علی تھیلے میں سے بلا ترتیب ایک گیند نکالتا ہے۔ احتمال معلوم کیجیے کہ نکالی گئی گیند۔ (الف) سفید ہو، (ب) سرخ ہو۔ (ج) سفید نہ ہو۔ (د) سیاہ نہ ہو۔}
\uq{408}{ایک عدد کو \lr{1} سے لے کر \lr{15} تک مکمل اعداد کے سیٹ سے ایک عدد بلا ترتیب منتخب کیا گیا۔ احتمال معلوم کیجیے کہ منتخب عدد ہو: (الف) طاق عدد (ب) \lr{3} کا جزو ضربی (ج) \lr{2} کا مربع}
\uq{409}{ایک عدد کو \lr{1} سے لے کر \lr{15} تک مکمل اعداد کے سیٹ سے ایک عدد بلا ترتیب منتخب کیا گیا۔ احتمال معلوم کیجیے کہ منتخب عدد ہو: (الف) مفرد عدد (ب) \lr{20}}
\uq{410}{ایک ڈائس کو دوبار پھینکا گیا۔ ہر بار پھینکنے پر \lr{4} سے بڑا عدد حاصل کرنے کا احتمال معلوم کیجیے۔}
\uq{411}{ایک بیگ میں \lr{5} نیلی گیندیں اور \lr{4} سبز گیندیں ہیں۔ ابو بکر تھیلے میں سے بلا ترتیب ایک گیند نکالتا ہے۔ احتمال معلوم کریں کہ نکالی گئی گیند \lr{(i)} نیلی ہے \lr{(ii)} نیلی نہیں ہے۔}
\uq{412}{ایک ساتھ ایک ڈائس اور ایک سکے کو پھینکنا جائے تو احتمال معلوم کریں کہ ڈائس پر طاق عدد اور سکہ پر ہیڈ آئے۔}
\uq{413}{مکانیت کا ایک کھیل جو کہ تیر کو گھمانے پر مشتمل ہے۔ جو کہ \lr{1،2،3،....،12} میں سے کسی ایک عدد پر رکنے کا یکساں امکان رکھتا ہے۔ احتمال کیا ہے کہ یہ تیر درج ذیل پر رکے گا۔ \lr{(i) 7} \lr{(ii)} ایک مفرد عدد \lr{(iii)} ایک مرکب عدد}
\uq{414}{دو ڈائس پھینکے جاتے ہیں۔ احتمال معلوم کریں کہ نتائج کا مجموعہ ہے: \lr{(i) 4} کے برابر \lr{(ii) 10} سے زیادہ \lr{(iii) 13} سے کم۔}
\uq{415}{اگر \lr{A} اور \lr{B} دو ایسے واقعات ہوں کہ \lr{$P(B)=\dfrac{1}{2},P(A)=\dfrac{1}{4}$} اور \lr{$P(A\cap B)=\dfrac{1}{8}$} ہوں تو \lr{$P(A\cup B)$} معلوم کیجیے۔}
\uq{416}{ایک اپارٹمنٹ میں گھر کے دروازوں کے نمبر \lr{1} سے \lr{50} تک ہیں۔ ان میں سے ایک گھر کا دروازہ نمبر حقیقت جفت عدد یا مکمل مربع ہونے کا احتمال معلوم کیجیے۔}
\uq{417}{ایک ٹیم کے کسی میچ کو جیتنے کا احتمال \lr{$\dfrac{3}{10}$} ہے اور کسی میچ کو ہارنے کا احتمال \lr{$\dfrac{2}{10}$} ہے۔ اس بات کا کیا احتمال ہے کہ \lr{(i)} ٹیم ایک خاص میچ جیت جائے یا ہار جائے۔ \lr{(ii)} ٹیم نہ ہی کوئی میچ جیتے نہ ہارے۔}
\uq{418}{دو ڈائس کے ایک ہی بار پھینکنے پر مجموعہ \lr{7} یا \lr{11} حاصل کرنے کا احتمال معلوم کیجیے۔}
\uq{419}{دو ڈائس کے ایک بار پھینکنے پر مجموعہ \lr{5} یا \lr{7} حاصل کرنے کا احتمال معلوم کیجیے۔}
\uq{420}{\lr{52} تاس کے پتوں سے بلا ترتیب ایک پتا نکالا گیا۔ نکالے گئے پتے کا احتمال معلوم کیجیے کہ وہ: \lr{(i)} ایک بادشاہ یا ایک غلام \lr{(ii)} نہ بادشاہ اور نہ ہی غلام ہو۔}
\uq{421}{ایک ڈائس کو دو بار پھینکا گیا۔ کیا احتمال ہے کہ دو میں سے کم از کم ایک بار پھینکنے پر ردد \lr{3} آئے؟}
\uq{422}{ایک تھیلے میں \lr{1} سے لے کر \lr{15} تک اعداد لکھے ہوئے \lr{15} کارڈ ہیں۔ ایک کارڈ بلا ترتیب نکالا گیا۔ احتمال معلوم کیجیے کہ اس پر \lr{5} کے جزو ضربی یا \lr{7} کے جزو ضربی والا عدد لکھا ہو۔}
\uq{423}{دو غیر جانبدار سکے اچھالے گئے۔ کم از کم ایک ہیڈ یا دو ہیڈ حاصل کرنے کا احتمال کیا ہے؟}
\uq{424}{ایک مصروف چوراہے پر \lr{50٪} گاڑیاں دائیں مڑتی ہیں۔ \lr{30٪} بائیں مڑتی ہیں اور \lr{20٪} سیدھی چلی جاتی ہیں۔ بلا ترتیب منتخب کی گئی گاڑی کے بائیں یا دائیں مڑنے کا احتمال کیا ہے؟}
\uq{425}{دو غیر جانبدار سکے اچھالے گئے۔ دو ہیڈ یا دو ٹیل حاصل کرنے کا کیا احتمال ہے؟}
\uq{426}{اگر \lr{A} اور \lr{B} ایک بلا ترتیب تجربے کے دو باہمی متناہی واقعات اور \lr{$P(\text{not }A)=0.45$} اور \lr{$P(A\cup B)=0.65$} ہو تو \lr{P(B)} معلوم کیجیے۔}
\uq{427}{اگر \lr{$P(A)=\dfrac{2}{3},P(B)=\dfrac{2}{5},P(A\cup B)=\dfrac{1}{3}$} ہو تو \lr{$P(A\cap B)$} معلوم کیجیے۔}
\uq{428}{\lr{A} اور \lr{B} دو ایسے واقعات ہیں کہ \lr{$P(A)=0.42,\ P(B)=0.48$} اور \lr{$P(A\cap B)=0.16$} تو معلوم کیجیے۔ \lr{(i)} \lr{$P(\overline{A})$} \lr{(ii)} \lr{$P(\overline{B})$} \lr{(iii)} \lr{$P(A\cup B)$}}
\uq{429}{پھولوں کی ایک ٹوکری میں \lr{50} پیلے \lr{70} سرخ اور \lr{80} سفید پھول ہیں۔ ان میں سے ایک پھول بلا ترتیب منتخب کیا جاتا ہے۔ پیلے یا سرخ پھول کے منتخب ہونے کا احتمال معلوم کیجیے۔}
\uq{430}{دو ڈائس ایک ساتھ پھینکے گئے۔ دونوں ڈائس پر ایک جیسے ہی اعداد ہوں یا اعداد کا مجموعہ \lr{10} حاصل کرنے کا احتمال معلوم کیجیے۔}
\uq{431}{عرشیہ منگل اور بدھ میں سے ایک دن اور صبح، شام میں سے ایک وقت منتخب کرتی ہے۔ امکانی خاکہ بنایئے۔ شجری خاکہ بھی بنایئے۔}
\uq{432}{\lr{1} سے \lr{10} کے اعداد میں سے دو اعداد کو واپس رکھنے کے ساتھ بلا ترتیب منتخب کیا گیا ہے۔ احتمال معلوم کیجیے کہ: \lr{(i)} دونوں مفرد ہوں \lr{(ii)} ان کا حاصل ضرب جفت ہو۔}
\uq{433}{لفظ \lr{PUNJAB} سے ایک حرف اور لفظ \lr{LAHORE} سے ایک اور حرف منتخب کیا گیا ہے۔ احتمال کیا ہے کہ: \lr{(i)} دونوں حروف واول ہوں \lr{(ii)} ایک غرف کانسوئٹ اور دوسرا اواؤول ہو۔}
\uq{434}{ایک ڈائس کو دو بار اچھالا گیا۔ احتمال معلوم کیجیے کہ ایک اچھال \lr{3} کا جزو ضربی اور دوسرا \lr{5} ہو۔}
\uq{435}{دو ڈائس پھنکی گئیں۔ ایک پر طاق اور دوسرے پر \lr{2} کا ضربی حاصل کرنے کا احتمال معلوم کیجیے۔}
\uq{436}{تاش کے پتوں میں سے دو پتے بلا ترتیب یکے بعد دیگرے بغیر واپسی کے نکالے گئے۔ احتمال معلوم کیجیے کہ پہلا پتا پان ہو اور دوسرا پتا بادشاہ ہو۔}
\uq{437}{تاش کے \lr{32} پتوں سے بغیر واپسی کے دو پتے منتخب کیے گئے۔ احتمال معلوم کیجیے کہ: \lr{(i)} دونوں سیاہ ہوں۔ \lr{(ii)} دو ملکہ ہوں \lr{(iii)} دونوں حکم ہوں۔ \lr{(iv)} دونوں اینٹیں ہوں۔}
\uq{438}{سلیم اچھی طرح پھینٹے ہوئے تاش کے \lr{52} پتوں میں سے بغیر واپسی کے یکے بعد دیگرے دو پتے نکالتا ہے۔ احتمال کیا ہے کہ پہلا پتا غلام ہو اور دوسرا پتا ملکہ ہے۔}
\uq{439}{تاش کے پتوں میں سے ایک پتا نکالا گیا ہے۔ ایک اکا اور ایک حکم حاصل کرنے کا احتمال معلوم کیجیے۔}
\uq{440}{فاطمہ تاش کے پتوں میں سے یکے بعد دیگرے (بغیر واپسی کے) دو پتے نکالتی ہے۔ احتمال معلوم کیجیے کہ دونوں پتے سرخ ہوں۔}
\uq{441}{اگر دو ڈائس پھینکے جائیں تو ان کے اوپری رخ اعاد کی حاصل ضرب \lr{6} یا اوپری رخ اعداد کا فرق \lr{5} حاصل کرنے کا احتمال معلوم کیجیے۔}
\uq{442}{ایک تھیلے میں \lr{7} نارنجی اور \lr{5} جامنی گولیاں ہیں۔ دو گولیاں یکے بعد دیگرے بغیر واپسی کے نکالی گئیں۔ احتمال معلوم کیجیے کہ پہلی گولی نارنجی اور دوسری گولی بھی نارنجی ہو۔}
\uq{443}{ایک تھیلے میں \lr{4} سرخ اور \lr{3} سبز گیندیں ہیں۔ کبیر ایک کے بعد ایک کر کے دو گیندیں نکالتا ہے (پہلی گیند واپس رکھنے کے بعد دوسری نکالی جاتی ہے)۔ احتمال معلوم کیجیے کہ پہلی گیند سرخ اور دوسری گیند بھی سرخ ہو۔}
\uq{444}{ایک تھیلے میں \lr{3} سفید اور \lr{2} سیاہ گولیاں ہیں۔ واپس رکھے بغیر یکے بعد دیگرے دو گولیاں نکالی جاتی ہیں۔ احتمال معلوم کیجیے کہ پہلی گولی، سفید ہو اور دوسری گولی بھی سفید ہو}
\uq{445}{ماجد \lr{52} تاش کے پتوں سے مسلسل (بغیر واپسی کے) دو پتے نکالتا ہے۔ احتمال معلوم کیجیے کہ پہلا پتا بادشاہ ہو اور دوسرا پتا ملکہ ہو۔}

\end{document}
